%% 第31回編集委員会（物理第5巻）：組 原W5
%% p5-09.tex  原-9　一方通行の電気：pn接合とLED・太陽電池（★1〜5・3段）
%% 中身は計画書 tools/phys5-plan.md の B「原-9」。公式の持ち主は A（A1・A40・A58〜A61）。
%%   ほかの組の章は、計画書 B の節の番号で指した（原-2 ③⑤、原-4 ②、原-6、原-8 ⑤⑦⑧、原-10 ③④、原-14 ②）。
%% 数値はすべて Python で計算した：checks/物理第5巻/原-9/calc_p5_09.py（出力 calc_p5_09_out.txt）、
%%   code_p5_09.py（⑫の囲みのプログラム）、kanji_p5.py（常用漢字の確かめ）。
%% 本文の丸め：k_B T＝0.0259 eV（300 K）、hc＝1240 eV·nm、e＝1.60×10⁻¹⁹ C、ε0＝8.85×10⁻¹² F/m（計画書 B のはじめ・E の8）。
%%   このため、ダイオードの「10倍で何 V」は 2.303×25.9＝59.6 mV（数学 3-12 の「約 0.06 V」と同じ。計画書の 59.5 は k_B T＝0.02585 で計算した値）。
%% 前の巻・数学の本との約束（計画書 A-6・C-6-2）
%%   ・A1（小学校：光電池）→ ①（ph:原-9:sec1）・例題11
%%   ・A40（高校物理：半導体・ダイオードの整流）→ ②⑦（ph:原-9:sec7）
%%   ・Q46（電-4:325 サーミスター、半導体は温度で抵抗率が下がる）→ ③・例題2（300 K→350 K で約28倍）
%%   ・Q8（熱-5:598 ペルチェ素子も第二法則の上限にしたがう）・Q47（電-4:581 熱電対の電圧を電子から）→ ⑬・例題13（COP の上限 9.26… を切り捨てて 9.2）
%%   ・Q49（電-4:678 ダイオード・トランジスタは半導体の性質から）→ ②⑥⑦（トランジスタは原-10）
%%   ・Q55（電-8:246 太陽光発電は光から直接）→ ⑩
%%   ・Q60（電-16:366「一方の向きにだけ電流を通す弁」）→ ⑦
%%   ・Q16（熱-15 ⑨ フェルミ分布で半導体の電子の数）→ ⑪（熱-15 の ε・μ と、この章の E・E_F は同じもの）
%%   ・M15（数学 3-12:714 ダイオードは約 0.06 V ごとに電流10倍）→ ⑦・例題8
%% 記号（この章）：電子と正孔の密度 n・p（この章で運動量は記号にしない）、真性キャリア密度 n_i、ドナー・アクセプターの密度 N_d・N_a
%%   （アボガドロ定数 N_A とぶつけない：N-4）、伝導帯の下のはし・価電子帯の上のはし E_c・E_v、禁制帯幅 E_g、フェルミ準位 E_F、
%%   有効状態密度 N_c・N_v、移動度 μ_n・μ_p、拡散係数 D_n・D_p、寿命 τ_n・τ_p、拡散距離 L_n・L_p、電場 \mathcal E（この章だけ。N-4。
%%   エネルギー E と分ける。第4巻の起電力 \mathcal E とは別の巻で、この章に起電力は出さない）、電位 φ、電荷密度 ρ_q（電-2 と同じ）、
%%   抵抗率 ρ、導電率 σ、内蔵電位 V_bi、空乏層の幅 W、飽和電流 I_s、光電流 I_L、電流密度 j（電-4 と同じ）、開放電圧 V_oc、曲線因子 FF、
%%   光の振動数 f（E＝hf：N-5）、面積 S、ゼーベック係数 S_ab は使わず「1 K あたりの電圧」とことばで書いた。
%% 出典（11月11日に確かめた）
%%   ・2014年ノーベル物理学賞（赤崎・天野・中村「明るく省エネルギーな白色光源を可能にした、効率のよい青色発光ダイオードの発明」）：
%%       https://www.nobelprize.org/prizes/physics/2014/summary/（curl で読んだ）。
%%   ・1956年ノーベル物理学賞（ショックレー・バーディーン・ブラッテン「半導体の研究とトランジスタ効果の発見」）：
%%       https://www.nobelprize.org/prizes/physics/1956/summary/（curl で読んだ）。
%%   ・定数：CODATA 2022（NIST、計画書 D4〜D6。k_B は physics.nist.gov/cgi-bin/cuu/Value?k を curl で読んだ）。
%%   ・LED の効率の低下（ドループ）：J. Iveland, L. Martinelli, J. Peretti, J. S. Speck, C. Weisbuch, Phys. Rev. Lett. 110, 177406 (2013)、
%%       arXiv:1304.5469（要旨を curl で読んだ。InGaN/GaN の LED で、電流を増やしたときの効率の低下と同時に、オージェ過程の高いエネルギーの
%%       電子を見つけ、ドループの主な原因はオージェ過程と結論）。ほかの原因（電子のもれ・欠陥）をあげる研究もあり、本文は「議論が続いている」と書いた。
%%   ・ショックレーとクワイサーの上限：W. Shockley, H. J. Queisser, J. Appl. Phys. 32, 510 (1961)（禁制帯 1.1 eV で約30 %。書誌は計画書と同じ）。
%%   ・ショックレーの pn 接合の理論：W. Shockley, Bell Syst. Tech. J. 28, 435 (1949)。
%% 〔要確認〕（一次資料に届かず、二次資料2つ以上で照らし合わせた値。11月11日）
%%   ・ケイ素（300 K）：禁制帯 1.12 eV、N_c＝2.8×10¹⁹・N_v＝1.04×10¹⁹ cm⁻³、移動度 1350・480 cm²/(V·s)、比誘電率 11.7。
%%       S. M. Sze『Physics of Semiconductor Devices』（2nd ed., 1981）の表の値として、半導体の表（eesemi.com「Properties of Si, Ge, and GaAs at 300K」）
%%       と、ほかの大学の講義資料・計算サイトの表（globalsino.com・firgelliauto.com）で照らし合わせた（WebSearch の抜粋）。
%%       Ioffe 研究所の NSM（www.ioffe.ru）は作業環境から届かなかった。出典の組（11月11日）：www.ioffe.ru・ミラーの www-f9.ijs.si は
%%       今回も応答がなかった。ほかの届く一次資料（NIST・IAEA など）にもケイ素の物性の表はなく、〔要確認〕のまま。
%%   ・真性キャリア密度：本・資料により 6.7×10⁹〜1.45×10¹⁰ cm⁻³ とちがう（上の式の値 6.7×10⁹、ふつうの教科書の 1.0×10¹⁰、Sze 第2版の 1.45×10¹⁰）。
%%       例題では 1.0×10¹⁰ を基準にし、本文でちがいを断った（計画書 C-3 の「どれを使うか章の組が決める」）。
%%   ・GaAs 1.42 eV、GaN 約3.4 eV、赤の LED（GaAsP・AlGaInP）約1.9 eV、青（InGaN）約2.7 eV：上の表と講義資料の値。本文では「約」とした。
%%   ・市販の太陽電池パネルの効率「20 % 前後〜25 % ほど」：pv magazine（2024年10月24日）の記事（Fraunhofer ISE が確かめた 25.4 % の記録、
%%       市販品 24.8 %）と NREL の Champion Module Efficiencies のページの検索の抜粋。20 % 前後は確かめられず〔要確認〕。
%%   ・ペルチェ素子の実際の成績係数（1 より小さいことが多い：温度差 20〜40 K で 0.3〜0.7 ほど）：検索の抜粋（Ferrotec の技術資料のページ、
%%       学位論文など）。一次資料は届かなかった〔要確認〕。
%%   ・例題の仮の値：ダイオードの I_s＝10⁻¹⁴ A（計画書 C-4 の約束。原-10 と同じ）、太陽電池の 35 mA/cm²・10⁻¹² A/cm²、光電池1個 0.6 V、
%%       LED の 2 V、寿命 1 μs、ペルチェの 1 対 0.4 mV/K は「〜とします」と書いた。
%%   ・やってみようの表は作ったデータ（6-4-7：光電流は光の強さに比例）。
\makeatletter
% 回路の部品（tikzpicture[x=1mm, y=1mm] の中。導線を先にかき、その上に置く。角は導線の向き）。第4巻 p4-09 と同じ形
\newcommand\txjR[2]{\begin{scope}[shift={#1}, rotate=#2]
  \filldraw[fill=white, draw=bkInk, line width=0.7pt] (-4.5,-1.6) rectangle (4.5,1.6);\end{scope}}
\newcommand\txjE[2]{\begin{scope}[shift={#1}, rotate=#2]
  \fill[white] (-1.2,-4.2) rectangle (1.2,4.2);
  \draw[bkInk, line width=0.7pt] (-1.0,-4.0) -- (-1.0,4.0);
  \draw[bkInk, line width=1.8pt] (1.0,-2.0) -- (1.0,2.0);\end{scope}}
% ダイオード（三角の向きに電流を通す）
\newcommand\txjD[2]{\begin{scope}[shift={#1}, rotate=#2]
  \fill[white] (-2.4,-2.6) rectangle (2.4,2.6);
  \filldraw[fill=bkInk!15, draw=bkInk, line width=0.7pt] (-2,-2.2) -- (-2,2.2) -- (2,0) -- cycle;
  \draw[bkInk, line width=1.1pt] (2,-2.2) -- (2,2.2);\end{scope}}
% LED（ダイオードに外向きの光の矢印）
\newcommand\txjLED[2]{\txjD{#1}{#2}\begin{scope}[shift={#1}, rotate=#2]
  \draw[bkOrange!85!black, line width=0.6pt, -{Stealth[length=1.4mm]}] (0.5,2.6) -- (2.2,5);
  \draw[bkOrange!85!black, line width=0.6pt, -{Stealth[length=1.4mm]}] (1.9,2.2) -- (3.6,4.6);\end{scope}}
% 結晶の原子（丸と記号）  \txjatom{(x,y)}{色}{字}
\newcommand\txjatom[3]{\filldraw[fill=#2, draw=bkInk, line width=0.5pt] #1 circle (2.6); \node[font=\sffamily\bfseries\footnotesize, text=bkInk] at #1 {#3};}
% 電子（黒い丸に −）と正孔（白い丸に ＋）
\newcommand\txjel[1]{\filldraw[fill=bkBlue!80!black, draw=bkInk, line width=0.4pt] #1 circle (1.1); \node[font=\bfseries\tiny, text=white, inner sep=0pt] at #1 {$-$};}
\newcommand\txjho[1]{\filldraw[fill=white, draw=bkRed!80!black, line width=0.7pt] #1 circle (1.1); \node[font=\bfseries\tiny, text=bkRed!80!black, inner sep=0pt] at #1 {$+$};}
\tikzset{txj lab/.style={font=\sffamily\footnotesize, text=bkInk}, txj arr/.style={line width=1pt, -{Stealth[length=2.2mm]}},
  txj dim/.style={bkGray, line width=0.5pt, {Stealth[length=1.8mm]}-{Stealth[length=1.8mm]}},
  txj band/.style={line width=1.3pt, bkInk}}
\makeatother

\phchapter{9}{一方通行の電気：\\pn接合とLED・太陽電池}{LEDは、どうやって電気を直接、光に変える？}{\bklevelrange{1}{5}}{一方通行の電気：pn接合とLED・太陽電池}
\bkmichishirube{3-12,5-14,5-7,6-4,8-13,熱-5,熱-15,電-2,電-4,電-16,原-2,原-4,原-8}{原-10,原-14}

\begin{goal}
\begin{itemize}[label={\color{bkTheme}●}, leftmargin=1.6\zw, itemsep=0pt, topsep=1pt]
\item \bklevel{1}\enspace 光電池に光を当てると電気ができ、LED は決まった向きにつないだときだけ光ることを、実験で確かめられる。
\item \bklevel{3}\enspace n形・p形の半導体の電子と正孔の数と抵抗率を求め、pn 接合が電流を一方の向きにだけ通すしくみを、図で説明できる。
\item \bklevel{3}\enspace ダイオードの電流と電圧の関係（約 60 mV で電流10倍）、LED の色と禁制帯の幅の関係を使える。
\item \bklevel{3}\enspace 太陽電池の電流と電圧のグラフから、開放電圧・いちばん大きな電力・曲線因子・効率を求められる。
\item \bklevel{5}\enspace フェルミ分布の裾からキャリアの数の式を、ドリフトと拡散のつり合いから内蔵電位とダイオードの式を導ける。
\end{itemize}
\end{goal}

\phkoushiki{%
\phk{半導体（n形・p形）、ダイオードの整流作用}{高校物理}{ph:原-9:sec7}
}

\begin{hatena}
\noindent\begin{minipage}[c]{0.56\linewidth}\raggedright
信号機・テレビ・家の照明は、この20年ほどで、ほとんど LED（発光ダイオード）に変わりました。白熱電球は、フィラメントを \qty{2000}{\celsius} 以上に熱して光らせるので、さわるととても熱くなります。ところが LED は、明るく光っていても、さわれないほど熱くはなりません。

LED は、電気を、熱を通らずに直接、光に変えているのです。いったい、どうやって？　そして、赤・緑・青の LED の色は、何で決まっているのでしょう。
\end{minipage}\hfill
\begin{minipage}[c]{0.40\linewidth}\centering
\begin{tikzpicture}[x=1mm, y=1mm]
  % 白熱電球
  \draw[bkInk, line width=0.8pt, fill=bkYellow!25] (10,22) circle (8);
  \fill[bkGray!40] (6,8) rectangle (14,14.5);
  \draw[bkInk, line width=0.6pt] (6,8) rectangle (14,14.5);
  \draw[bkOrange!85!black, line width=0.9pt] (7.5,22) -- (8.5,24) -- (9.5,20) -- (10.5,24) -- (11.5,20) -- (12.5,22);
  \foreach \a in {20,50,80,110,140,170} \draw[bkRed!70!black, line width=0.6pt, decorate, decoration={snake, amplitude=0.5mm, segment length=2mm}] (10,22) ++(\a:9.5) -- ++(\a:4);
  \node[txj lab, anchor=north] at (10,6) {白熱電球};
  \node[txj lab, text=bkRed!80!black, anchor=north] at (10,2.3) {熱い};
  % LED
  \draw[bkInk, line width=0.8pt, fill=bkBlue!12] (34,14) -- (34,24) arc (180:0:4) -- (42,14) -- cycle;
  \draw[bkInk, line width=0.7pt] (36,14) -- (36,6) (40,14) -- (40,4);
  \fill[bkBlue!60!black] (37,17) rectangle (39,19);
  \foreach \a in {50,90,130} \draw[bkBlue!70!black, line width=0.8pt, -{Stealth[length=1.6mm]}] (38,22) ++(\a:7) -- ++(\a:5);
  \node[txj lab, anchor=north] at (38,2.3) {LED};
  \node[txj lab, text=bkBlue!70!black, anchor=north] at (38,-1.4) {熱くない};
\end{tikzpicture}
\end{minipage}
\end{hatena}

\noindent 電気を一方の向きにだけ通す部品を\textbf{ダイオード}といいます。\bkidx{だいおーど}{ダイオード}物理シリーズの『電磁気学』では、ダイオードを「一方の向きにだけ電流を通す弁」、トランジスタを「オンかオフだけのスイッチ」として、中身を見ずに使いました（電-12・電-16）。この章では、その中身を開けます。中身は、ケイ素などの\textbf{半導体}に、ほんの少しだけ別の原子を混ぜたものを、2種類はり合わせたもの（\textbf{pn 接合}）です。LED・太陽電池・カメラの光センサーも、同じ pn 接合でできています。前の章（原-8）で学んだ「エネルギーの帯（バンド）」が、ここで道具になります。

\begin{phmijika}[身の回りの pn 接合]
\noindent\begin{tikzpicture}[x=1mm, y=1mm]
  \foreach \i/\a/\b in {0/{LED の照明}/{電気→光}, 1/{太陽光パネル}/{光→電気}, 2/{スマホのカメラ}/{光→信号}, 3/{充電器}/{交流→直流}} {
    \begin{scope}[shift={(\i*30,0)}]
      \draw[bkInk, line width=0.6pt, rounded corners=1.5mm, fill=bkTheme!6] (0,0) rectangle (26,20);
      \node[txj lab, anchor=north] at (13,-1) {\a};
      \node[txj lab, text=bkGray!80!black, anchor=north] at (13,-5.2) {\b};
    \end{scope}}
  % LED の照明
  \draw[bkInk, line width=0.7pt, fill=bkYellow!30] (8,6) -- (8,12) arc (180:0:5) -- (18,6) -- cycle;
  \fill[bkGray!50] (9.5,3) rectangle (16.5,6);
  \foreach \a in {40,90,140} {\draw[bkOrange!85!black, line width=0.6pt] (13,13) ++(\a:6.2) -- ++(\a:2.2);}
  % 太陽光パネル
  \begin{scope}[shift={(30,0)}]
    \fill[bkBlue!40] (4,5) -- (20,5) -- (23,14) -- (7,14) -- cycle;
    \foreach \t in {0.25,0.5,0.75} {\draw[white, line width=0.5pt] ({4+16*\t},5) -- ({7+16*\t},14);}
    \draw[white, line width=0.5pt] (5.5,9.5) -- (21.5,9.5);
    \draw[bkInk, line width=0.6pt] (4,5) -- (20,5) -- (23,14) -- (7,14) -- cycle;
    \fill[bkOrange!80] (22,18) circle (1.6);
  \end{scope}
  % スマホのカメラ
  \begin{scope}[shift={(60,0)}]
    \draw[bkInk, line width=0.7pt, rounded corners=1.2mm, fill=bkGray!15] (8,3) rectangle (18,18);
    \draw[bkInk, line width=0.6pt, fill=bkNavy!60] (11.5,14.5) circle (1.6);
    \draw[bkInk, line width=0.6pt, fill=bkNavy!60] (14.8,14.5) circle (1.2);
  \end{scope}
  % 充電器
  \begin{scope}[shift={(90,0)}]
    \draw[bkInk, line width=0.7pt, rounded corners=1mm, fill=white] (7,5) rectangle (17,15);
    \draw[bkInk, line width=0.9pt] (9.5,15) -- (9.5,18.5) (14.5,15) -- (14.5,18.5);
    \draw[bkInk, line width=0.7pt] (12,5) -- (12,2.5) -- (21,2.5);
  \end{scope}
\end{tikzpicture}\par\vspace{2pt}
4つとも、中に pn 接合があります。はたらきは、電気から光、光から電気、光から信号、交流から直流と、違って見えます。けれども、どれも「n形と p形の境目にできる坂」を、電子と正孔がのぼるか下りるかの話です。この章では、その坂を1つずつ調べます。
\end{phmijika}

\begin{yattemiyou}[LED と光電池で、一方通行を確かめよう]
\textbf{用意するもの}\par
乾電池（単3、2本）と電池ボックス、抵抗（\qty{100}{\ohm}）、赤と青の LED（ふつうの砲弾形）、ミノムシクリップつきの導線、理科室の光電池（モーターつきのもの）、テスター\par\vspace{2pt}
\noindent\begin{minipage}[c]{0.64\linewidth}
\begin{enumerate}[label={\bkcircled[bkGreen]{\arabic*}}, leftmargin=2.4\zw, itemsep=1pt, topsep=0pt]
\item 乾電池2本（\qty{3}{V}）・抵抗・赤い LED を、図のように1つの輪につなぐ。LED の足の長いほうを、電池の＋の側にする。
\item LED の向きだけを反対にして、光るかどうかを見る。青い LED でも、①②をする。
\item 光電池をモーターにつなぎ、手でかげを作ったり、光電池を光の向きに傾けたりして、回り方を見る。
\item 光電池をテスター（電流の、いちばん小さいはんい）につなぎ、光電池の半分を紙でかくすと、電流がどうなるかを調べる。
\end{enumerate}
\end{minipage}\hfill
\begin{minipage}[c]{0.33\linewidth}\centering
\begin{tikzpicture}[x=1mm, y=1mm]
  \draw[bkInk, line width=0.7pt] (8,0) -- (8,24) -- (34,24) -- (34,0) -- cycle;
  \txjE{(8,12)}{-90}
  \txjR{(21,24)}{0}
  \txjLED{(21,0)}{180}
  \node[txj lab, anchor=east] at (5,12) {3 V};
  \node[txj lab, anchor=south] at (21,26.5) {100 Ω};
  \node[txj lab, anchor=north west] at (24,-2.2) {LED};
  \draw[bkInk, txj arr] (37,18) -- (37,6);
\end{tikzpicture}
\end{minipage}
\par\vspace{6pt}\textbf{気をつけること}：家のコンセントにはつながない。LED は、抵抗をつながずに電池に直接つなぐと、大きな電流が流れてこわれることがある。乾電池を3本以上つながない。
\end{yattemiyou}

\begin{bkcenter}
\begin{bktable}{colspec={Q[l,m,wd=58mm] Q[l,m,wd=62mm]}, rowsep=1pt, row{1}={bg=bkTheme!10}, cells={font=\small}}
したこと & 結果 \\
赤い LED、足の長いほうを＋の側に & 明るく光る \\
赤い LED、向きを反対に & 光らない \\
青い LED、足の長いほうを＋の側に & 光るが、赤より暗い \\
光電池（全面に日光） & モーターがよく回る。テスターで約 \qty{60}{mA} \\
光電池を半分かくす & モーターがゆっくり回る。約 \qty{30}{mA} \\
\end{bktable}
\bktabcaption{やってみようで、ソラさんが調べた結果}\label{ph:原-9:tab-yattemiyou}
\end{bkcenter}

\noindent やってみようの答え：表\ref{ph:原-9:tab-yattemiyou}から、次のことがわかります。
\begin{enumerate}[label={\bkcircled{\arabic*}}, leftmargin=2.4\zw, itemsep=1pt, topsep=2pt]
\item LED は、決まった向きにつないだときだけ光る。反対にすると、電流が流れない（一方通行）。
\item 同じ電池と抵抗でも、青い LED は赤より暗い。青い LED は、光らせるのに、赤より大きな電圧が要るからです（⑨）。
\item 光電池は、光が当たると、電池のように電流を流す。光が当たる面積が半分になると、電流も半分になる。
\end{enumerate}
①は⑦のダイオードのしくみ、②は⑨の LED の色、③は⑩の太陽電池で説明します。LED と光電池は、じつは同じ「pn 接合」で、電気から光を出すか、光から電気を出すかの向きが違うだけなのです。

\begin{omoidasou}
\bkrecall{原-8}{原子がたくさん並んだ固体では、電子のエネルギーは帯（バンド）になる。電子がつまった\textbf{価電子帯}と、空いた\textbf{伝導帯}の間には、電子が入れない\textbf{禁制帯}がある。そのはばが\textbf{禁制帯幅} $E_{\mathrm g}$。ケイ素は約 \qty{1.12}{eV}、ガラスは約 \qty{9}{eV}。}
\bkrecall{原-2}{光はエネルギー $E=hf$ のつぶ（光子）。波長 $\lambda$ の光子のエネルギーは $E=hc/\lambda$。$hc=\qty{1240}{eV.nm}$ なので、$E$〔eV〕$=1240/\lambda$〔nm〕。}
\bkrecall{熱-15}{温度 $T$ で、エネルギーが $\Delta E$ だけ高い状態にいる割合は、ボルツマン因子 $e^{-\Delta E/(k_{\mathrm B}T)}$ に比例する。300 K で $k_{\mathrm B}T=\qty{0.0259}{eV}$。}
\bkrecall{電-4}{電流の密度は $j=\sigma\mathcal E$（$\sigma$ は導電率、抵抗率 $\rho=1/\sigma$）。電子の数密度を $n$、平均の速さを $v$ とすると、$j=env$。}
\bkrecall{3-12}{$y=10^x$ は、$x$ が1増えるごとに $y$ が10倍。片対数グラフでは、指数関数は直線になる。$\ln10\approx2.303$。}
\end{omoidasou}

\noindent この章の記号：電子の数密度を $n$、正孔（②で説明）の数密度を $p$ と書きます。この章では、運動量を記号にしません。半導体の中の電場は、エネルギー $E$ と区別するため、$\mathcal E$（イー、筆記体）と書きます。第4巻の起電力 $\mathcal E$ とは別のものです。禁制帯幅は $E_{\mathrm g}$ です（原-8）。

%% ======================== 物理の道具 ========================
\Needspace{18\baselineskip}\bksection{doguu}{物理の道具}

\phstage{me}{光電池と LED で、半導体にふれる}{1}

\bksubsection{1}{光電池で電気をつくる}
\label{ph:原-9:sec1}
小学校の理科で使った光電池は、ケイ素（シリコン）の板でできています。光が当たっているあいだ、光電池は電池のように、つないだモーターや豆電球に電流を流します。光を強くするほど、電流は大きくなります。光電池を2枚並べて（並列に）同じ光を当てると、電流は2倍になります。
\begin{phbaisuru}{光の強さ}
\noindent\begin{minipage}[c]{0.48\linewidth}\raggedright
光電池にあてる光の強さを2倍にすると、流れる電流（ショートさせたときの電流）も2倍になります。電流が光の強さに比例するのは、光のつぶ（光子）1個が、電子を1個ずつ押し出しているからです（⑩）。光を2倍にすると、1秒に当たる光子の数が2倍になり、押し出される電子の数も2倍になるのです。表\ref{ph:原-9:tab-yattemiyou}の「半分かくすと半分」も、同じ理由です。
\end{minipage}\hfill
\begin{minipage}[c]{0.48\linewidth}\centering
\begin{tikzpicture}
\begin{axis}[bkplot, axis lines=left, width=48mm, height=40mm, xmin=0, xmax=2.3, ymin=0, ymax=140, xtick={0,1,2}, ytick={0,60,120}, xlabel={光の強さ}, ylabel={電流（mA）}]
  \addplot[bkcurve, domain=0:2.15] {60*x};
  \fill[bkInk] (axis cs:0.5,30) circle (1.6pt) (axis cs:1,60) circle (1.6pt) (axis cs:2,120) circle (1.6pt);
\end{axis}
\end{tikzpicture}
\end{minipage}
\end{phbaisuru}
ふしぎなのは、光電池の中には、電池のような薬品も、発電機のような回る部品もないことです。ケイ素の板の中で、光が直接、電気のエネルギーに変わっています。そのしくみを知るために、まず、ケイ素の中の電子を見ていきましょう。

\Needspace{8\baselineskip}\bksubsection{2}{半導体の電子と正孔：n形とp形}
\label{ph:原-9:sec2}
ケイ素の原子は、いちばん外側に4個の電子をもっています。結晶の中では、1つの原子が4つのとなりの原子と、電子を1個ずつ出し合って、2個で1組の「結合の手」を作っています（図\ref{ph:原-9:fig-kessho}の左）。本当は立体（ダイヤモンドと同じ形。原-8 ①）ですが、ここでは平らにかきます。

\begin{bkfigure}
\begin{tikzpicture}[x=1mm, y=1mm]
  \foreach \ox/\ttl/\kind in {0/{まじりけのないケイ素}/0, 45/{n形（リンを入れる）}/1, 90/{p形（ホウ素を入れる）}/2} {
    \begin{scope}[shift={(\ox,0)}]
      \foreach \i in {0,1,2} \foreach \j in {0,1,2} {
        \draw[bkGray, line width=0.5pt] (6+12*\i-6,6+12*\j) -- (6+12*\i+6,6+12*\j);
        \draw[bkGray, line width=0.5pt] (6+12*\i,6+12*\j-6) -- (6+12*\i,6+12*\j+6);}
      \foreach \i in {0,1,2} \foreach \j in {0,1,2} {
        \fill[bkInk] (6+12*\i-4.4,6+12*\j+0.9) circle (0.55) (6+12*\i-4.4,6+12*\j-0.9) circle (0.55);
        \fill[bkInk] (6+12*\i+0.9,6+12*\j+4.4) circle (0.55) (6+12*\i-0.9,6+12*\j+4.4) circle (0.55);}
      \foreach \i in {0,1,2} \foreach \j in {0,1,2} {\txjatom{(6+12*\i,6+12*\j)}{bkGray!15}{Si}}
      \node[txj lab, anchor=north] at (18,-1.5) {\ttl};
    \end{scope}}
  % n形：まん中をリンに、余った電子
  \begin{scope}[shift={(45,0)}]
    \txjatom{(18,18)}{bkBlue!25}{P}
    \txjel{(23.6,22.6)}
    \draw[bkBlue!70!black, line width=0.6pt] (24.2,23.6) -- (27,38);
    \node[txj lab, text=bkBlue!70!black, anchor=south] at (27,38) {余った電子};
  \end{scope}
  % p形：まん中をホウ素に、1つ足りない（正孔）
  \begin{scope}[shift={(90,0)}]
    \txjatom{(18,18)}{bkRed!20}{B}
    \fill[white] (12.6,17.1) rectangle (14.6,20.0);
    \txjho{(13.6,18.9)}
    \node[txj lab, text=bkRed!80!black, anchor=south] at (12,38) {正孔（電子の抜けた穴）};
    \draw[bkRed!70!black, line width=0.6pt] (12,38) -- (13.4,20);
  \end{scope}
\end{tikzpicture}
\bkfigcaption{ケイ素の結晶を平らにかいた図（丸は原子、線にそった2つの点は結合の手の電子）。左：手がすべて電子でうまっている。中：電子を5個もつリン（P）が入ると、電子が1個余る。右：電子を3個しかもたないホウ素（B）が入ると、手が1か所欠ける（白い丸の＋）。}\label{ph:原-9:fig-kessho}
\end{bkfigure}

\phboxhead{まじりけのないケイ素は、電気を通しにくい}
結合の手が全部うまっていると、電子は、手から手へ動けません。動くには、結合の手から外れて、結晶の中を自由に動ける状態になる必要があります。それに要るエネルギーが、原-8 の禁制帯幅 $E_{\mathrm g}=\qty{1.12}{eV}$ です。室温の熱のエネルギー（$k_{\mathrm B}T=\qty{0.0259}{eV}$）は、その約40分の1しかないので、外れる電子はごくわずかです。だから、まじりけのないケイ素は、金属よりずっと電気を通しにくいのです。

\phboxhead{正孔：電子の抜けた穴も、電気を運ぶ}
電子が1個、結合の手から外れると、その手には電子1個分の「穴」が残ります。穴の周りは、電子が1個足りないので、＋の電荷 $e$ をもつように見えます。となりの手の電子がこの穴に移ると、穴は反対の向きに1つ動きます。電場をかけると、電子は＋極のほうへ、穴は－極のほうへ動き、どちらも同じ向きの電流になります。この穴を、＋の電荷をもつ1つの粒子のように考えて\textbf{正孔}（ホール）とよびます。\bkidx{せいこう}{正孔}電子と正孔のように、電気を運ぶものをまとめて\textbf{キャリア}といいます。\bkidx{きゃりあ}{キャリア}

\begin{phmijika}[劇場の席の空き]
満席の劇場で、前のほうの席が1つ空いたとします。うしろの人が1つずつ前へつめていくと、空いた席は、うしろへうしろへと動いていきます。動いているのは人（電子）ですが、空席（正孔）が1つ、反対向きに動くと考えたほうが、ずっと簡単です。
\end{phmijika}

\phboxhead{n形とp形：混ぜる原子で、キャリアを選ぶ}
ケイ素に、電子を5個もつリン（P）を少しだけ混ぜます。リンは、4個の電子で結合の手を作り、1個が余ります（図\ref{ph:原-9:fig-kessho}の中）。この電子は、とても小さなエネルギー（約 \qty{0.045}{eV}）でリンから離れ、自由に動けるようになります。電子を与えるので、リンのような原子を\textbf{ドナー}といい、電子が主にキャリアになる半導体を\textbf{n形}（negative の n）といいます。\bkidx{えぬがたはんどうたい}{n形半導体}

反対に、電子を3個しかもたないホウ素（B）を混ぜると、結合の手が1か所欠け、正孔ができます（図\ref{ph:原-9:fig-kessho}の右）。電子を受けとるので\textbf{アクセプター}といい、正孔が主にキャリアになる半導体を\textbf{p形}（positive の p）といいます。\bkidx{ぴーがたはんどうたい}{p形半導体}どちらも、半導体全体では電荷は0です。リンの原子核の＋の電荷が1つ多いぶん、余った電子とつり合っているからです。

\begin{bkfigure}
\begin{tikzpicture}[x=1mm, y=1mm]
  \foreach \ox/\ttl in {0/{まじりけのない},42/{n形},84/{p形}} {
    \begin{scope}[shift={(\ox,0)}]
      \fill[bkBlue!10] (0,24) rectangle (32,32);
      \fill[bkRed!10] (0,0) rectangle (32,8);
      \draw[txj band] (0,24) -- (32,24);
      \draw[txj band] (0,8) -- (32,8);
      \node[txj lab, anchor=south] at (16,32.5) {\ttl};
    \end{scope}}
  \node[txj lab, anchor=east] at (-1,28) {伝導帯};
  \node[txj lab, anchor=east] at (-1,4) {価電子帯};
  \node[txj lab, anchor=east, align=right] at (-1,16) {禁制帯\\$E_{\mathrm g}$};
  % まじりけのない：同じ数
  \txjel{(10,26.5)} \txjel{(22,26.5)}
  \txjho{(10,5.5)} \txjho{(22,5.5)}
  % n形：ドナーの準位（破線）と、たくさんの電子
  \draw[bkBlue!70!black, line width=0.8pt, dashed] (42,21.5) -- (74,21.5);
  \node[txj lab, text=bkBlue!70!black, anchor=north] at (58,21) {ドナーの準位};
  \foreach \x in {45,50,55,60,65,70} {\txjel{(\x,26.5)}}
  \txjho{(58,5.5)}
  % p形：アクセプターの準位（点線）と、たくさんの正孔
  \draw[bkRed!70!black, line width=0.8pt, dotted] (84,10.5) -- (116,10.5);
  \node[txj lab, text=bkRed!70!black, anchor=south] at (100,11) {アクセプターの準位};
  \foreach \x in {87,92,97,102,107,112} {\txjho{(\x,5.5)}}
  \txjel{(100,26.5)}
\end{tikzpicture}
\bkfigcaption{帯の図でみた3つの半導体（上が電子のエネルギーの高い向き。黒い丸が電子、白い丸が正孔）。n形では、伝導帯のすぐ下（破線）のドナーから電子が伝導帯へ上がる。p形では、価電子帯のすぐ上（点線）のアクセプターが価電子帯から電子を受けとり、価電子帯に正孔が残る。}\label{ph:原-9:fig-band3}
\end{bkfigure}
帯の図（原-8）でかくと、図\ref{ph:原-9:fig-band3}のようになります。電子は伝導帯の中を、正孔は価電子帯の中を動きます。混ぜる量は、ケイ素の原子の数の100万分の1よりずっと少なくてよく、それだけで電気の通しやすさが何万倍も変わります（例題4）。このように、まじりけを入れることを\textbf{ドーピング}といいます。

\begin{phmijika}[電子体温計のサーミスター]
半導体は、温度が上がると、結合の手から外れる電子が増えて、電気を通しやすくなります（物理シリーズの『電磁気学』（電-4 ④））。電子体温計の先には、この性質を使った温度センサー（サーミスター）が入っています。どれくらい増えるのかを、次の③で数えます。
\end{phmijika}

\Needspace{12\baselineskip}\phstage{te}{電子と正孔を数えて、pn 接合をつくる}{3}

\bksubsection{3}{温度で増える電子と正孔：真性キャリア密度}
\label{ph:原-9:sec3}
まじりけのない半導体では、結合の手から電子が1個外れるたびに、電子と正孔が1組できます。だから、電子の数密度 $n$ と正孔の数密度 $p$ は等しく、これを\textbf{真性キャリア密度} $n_{\mathrm i}$ といいます（i は intrinsic＝そのものの）。\bkidx{しんせいきゃりあみつど}{真性キャリア密度}

外れるには、禁制帯をこえるエネルギー $E_{\mathrm g}$ が要ります。熱でたまたまそれだけのエネルギーをもらえる割合は、熱-15 のボルツマン因子のような形になります。詳しい計算（⑪）の結果は、
\pheq{n_{\mathrm i}=\sqrt{N_{\mathrm c}N_{\mathrm v}}\;e^{-E_{\mathrm g}/(2k_{\mathrm B}T)}}{cm$^{-3}$ ＝ cm$^{-3}$}
\label{ph:原-9:ni}です。$N_{\mathrm c}$・$N_{\mathrm v}$ は、伝導帯・価電子帯の「席の数」にあたる量（\textbf{有効状態密度}）で、ケイ素では 300 K で $N_{\mathrm c}=\qty{2.8e19}{cm^{-3}}$、$N_{\mathrm v}=\qty{1.04e19}{cm^{-3}}$ ほどです（本によって少し違います。温度の $\tfrac32$ 乗に比例してゆっくり変わります）。

\phboxhead{半導体の数密度の単位}
半導体の本では、数密度を 1 cm$^3$ あたりの個数（cm$^{-3}$）で書くのがふつうです。1 m$^3$＝$10^6$ cm$^3$ なので、$\qty{1}{cm^{-3}}=\qty{e6}{m^{-3}}$ です。たとえば $\qty{1.0e10}{cm^{-3}}=\qty{1.0e16}{m^{-3}}$ です。この章も cm$^{-3}$ で書きます。移動度も、本の慣用どおり cm$^2$/(V·s) で書きます。

\begin{phyomi}[$n_{\mathrm i}=\sqrt{N_{\mathrm c}N_{\mathrm v}}\,e^{-E_{\mathrm g}/(2k_{\mathrm B}T)}$]
「電子と正孔の数は、席の数 × 禁制帯の半分をこえる割合」と読みます。肩が $E_{\mathrm g}$ ではなく $E_{\mathrm g}/2$ なのは、電子と正孔が1組でできるので、2つの数のかけ算 $np$ が $e^{-E_{\mathrm g}/(k_{\mathrm B}T)}$ に比例し、その平方根をとるからです（④・⑪）。ケイ素の値を入れると $n_{\mathrm i}\approx\qty{7e9}{cm^{-3}}$ です（例題2）。ケイ素の原子は 1 cm$^3$ に $\num{5.0e22}$ 個あるので（原-8 ③）、原子 $10^{13}$ 個に1組しか、電子と正孔がありません。
\end{phyomi}

\noindent 本や資料によって、ケイ素の $n_{\mathrm i}$ は $\qty{0.7e10}{cm^{-3}}$ から $\qty{1.5e10}{cm^{-3}}$ ほどまで、違う値が書かれています。肩の $E_{\mathrm g}/(2k_{\mathrm B}T)\approx22$ が大きいので、$E_{\mathrm g}$ や $N_{\mathrm c}$・$N_{\mathrm v}$ のわずかな違いが、何十％もの違いになるからです。この章の例題では、多くの教科書と同じ $n_{\mathrm i}=\qty{1.0e10}{cm^{-3}}$ を使います。

\begin{phbaisuru}[300 K から 350 K]{温度}
\noindent\begin{minipage}[c]{0.45\linewidth}\raggedright
温度を 300 K（27 ℃）から 350 K（77 ℃）へ、たった17 \% 上げるだけで、$n_{\mathrm i}$ は約28倍になります（例題2）。肩の $E_{\mathrm g}/(2k_{\mathrm B}T)$ が 21.7 から 18.6 へ小さくなり、$e^{3.1}\approx22$ 倍、席の数の増え方 $(350/300)^{3/2}\approx1.26$ 倍をかけて28倍です。縦軸を対数にしても、50 K ごとに1けた以上も上がります。電-4 で見た「半導体は温度が上がると抵抗率が大きく下がる」のは、このためです。
\end{minipage}\hfill
\begin{minipage}[c]{0.52\linewidth}\centering
\begin{tikzpicture}
\begin{axis}[bkplot, axis lines=left, ymode=log, width=54mm, height=46mm, xmin=240, xmax=510, ymin=1e7, ymax=1e14, xtick={300,400,500}, ytick={1e8,1e10,1e12,1e14}, yticklabels={$10^8$,$10^{10}$,$10^{12}$,$10^{14}$}, xlabel={温度（K）}, ylabel={$n_{\mathrm i}$（cm$^{-3}$）}]
  \addplot[bkcurve, domain=245:505, samples=60] {sqrt(2.8e19*1.04e19*(x/300)^3)*exp(-1.12/(2*8.617333e-5*x))};
  \fill[bkInk] (axis cs:300,6.68e9) circle (1.6pt) (axis cs:350,1.86e11) circle (1.6pt);
  \node[font=\sffamily\footnotesize, anchor=west] at (axis cs:303,3.0e9) {300 K};
  \node[font=\sffamily\footnotesize, anchor=west] at (axis cs:355,8.0e10) {350 K};
\end{axis}
\end{tikzpicture}
\end{minipage}
\end{phbaisuru}

\Needspace{12\baselineskip}\bksubsection{4}{混ぜる原子で数を決める：$np=n_{\mathrm i}^{\,2}$}
\label{ph:原-9:sec4}
電子と正孔は、出会うと、電子が穴にはまって、両方とも消えます（\textbf{再結合}）。半導体の中では、熱で電子と正孔の組が生まれるのと、出会って消えるのとが、いつも同時に起きていて、つり合っています。

消える速さは、電子の数と正孔の数の両方に比例します（出会う回数は、両方が多いほど多い）。生まれる速さは、温度と $E_{\mathrm g}$ だけで決まります。だから、つり合いでは、まじりけがあってもなくても、
\pheq{np=n_{\mathrm i}^{\,2}}{cm$^{-6}$ ＝ cm$^{-6}$}
\label{ph:原-9:np}が成り立ちます。\bkidx{しつりょうさようのほうそく}{質量作用の法則}化学の「質量作用の法則」（化学平衡で、濃さのかけ算が一定）と同じ形です。

\begin{phyomi}[$np=n_{\mathrm i}^{\,2}$]
「電子と正孔の数のかけ算は、温度で決まる一定の値」と読みます。片方を増やすと、もう片方はその分だけ減ります。リンを $N_{\mathrm d}=\qty{1e16}{cm^{-3}}$ 入れると、電子は $n\approx N_{\mathrm d}=10^{16}$（$n_{\mathrm i}$ の100万倍）になり、正孔は $p=n_{\mathrm i}^{\,2}/n=10^4$（$n_{\mathrm i}$ の100万分の1）に減ります（例題3）。多いほうのキャリアを\textbf{多数キャリア}、少ないほうを\textbf{少数キャリア}といいます。\bkidx{しょうすうきゃりあ}{少数キャリア}
\end{phyomi}

\noindent n形で $n\approx N_{\mathrm d}$ としてよいのは、室温ではドナーのほとんど全部が電子を出しているからです（ドナーの準位は伝導帯のすぐ下、\qty{0.045}{eV} ほどしか離れていない）。そして、ドナーから出る電子 $10^{16}$ に比べて、熱で結合の手から外れる電子 $n_{\mathrm i}=10^{10}$ は無視できるほど少ないからです。p形も同じで、$p\approx N_{\mathrm a}$、$n=n_{\mathrm i}^{\,2}/N_{\mathrm a}$ です。

\Needspace{12\baselineskip}\bksubsection{5}{電子が動く2つのしかた：ドリフトと拡散}
\label{ph:原-9:sec5}
半導体の中を電流が流れるしかたは、2つあります。

\phboxhead{ドリフト：電場に押されて流れる}
電場 $\mathcal E$ をかけると、電子は電場と反対向きに、正孔は電場の向きに、平均の速さ $v$ で動きます。速さは電場に比例し、その比例定数を\textbf{移動度} $\mu$ といいます（$v=\mu\mathcal E$）。\bkidx{いどうど}{移動度}ケイ素（300 K）の電子は $\mu_{\mathrm n}$ が約 1350、正孔は $\mu_{\mathrm p}$ が約 480（単位 cm$^2$/(V·s)）です。電子と正孔の両方の流れを足すと、電-4 の $j=env$ から
\pheq{\sigma=e\,(n\mu_{\mathrm n}+p\mu_{\mathrm p})}{S/cm ＝ C × cm$^{-3}$ × cm$^2$/(V·s)}
\label{ph:原-9:sigma}が導電率です（抵抗率は $\rho=1/\sigma$）。キャリアの数を混ぜる量で何けたも変えられるので、半導体の抵抗率は、ドーピングで自由に決められます（例題4）。

\phboxhead{拡散：濃いほうから、うすいほうへ広がる}
電場がなくても、キャリアが多い所から少ない所へは、キャリアが広がっていきます。水にたらしたインクが広がるのと同じで、それぞれのキャリアは勝手な向きに動きますが、多い所から出ていく数のほうが、入ってくる数より多いからです。1秒に1 cm$^2$ を通る数は、数密度の坂（1 cm あたりの変わり方）に比例します。
\pheq{\text{拡散の流れ}=-D\,\frac{\dd n}{\dd x}}{cm$^{-2}$/s ＝ cm$^2$/s × cm$^{-4}$}
\label{ph:原-9:diff}比例定数 $D$ を\textbf{拡散係数}といいます。\bkidx{かくさんけいすう}{拡散係数}負の符号は「坂を下る向き」の意味です。これは、熱-11 で学んだフーリエの法則（熱は温度の坂を下る）と、まったく同じ形です。

\begin{bkfigure}
\begin{tikzpicture}[x=1mm, y=1mm]
  % 左：ドリフト
  \begin{scope}
    \draw[bkInk, line width=0.6pt] (0,0) rectangle (52,22);
    \foreach \x/\y in {6/5,13/16,19/8,26/14,32/4,38/17,45/9,48/15} {\txjel{(\x,\y)}}
    \foreach \x/\y in {6/5,19/8,32/4,45/9} {\draw[bkBlue!70!black, line width=0.6pt, -{Stealth[length=1.4mm]}] (\x-1.6,\y) -- ++(-4,0);}
    \draw[bkGreen!50!black, line width=1.1pt, -{Stealth[length=2.2mm]}] (8,26) -- (44,26);
    \node[txj lab, text=bkGreen!40!black, anchor=south] at (26,27) {電場 $\mathcal E$};
    \node[txj lab, anchor=north] at (26,-1.5) {ドリフト：電場と反対向きに動く};
  \end{scope}
  % 右：拡散
  \begin{scope}[shift={(62,0)}]
    \draw[bkInk, line width=0.6pt] (0,0) rectangle (52,22);
    \foreach \x/\y in {3/4,4/12,5/19,8/8,9/16,11/3,12/11,14/18,16/7,18/14,21/4,24/17,27/10,30/5,34/15,39/9,46/13} {\txjel{(\x,\y)}}
    \draw[bkBlue!70!black, line width=1.1pt, -{Stealth[length=2.2mm]}] (14,26) -- (42,26);
    \node[txj lab, text=bkBlue!70!black, anchor=south] at (28,27) {電子の流れ};
    \node[txj lab, anchor=north] at (26,-1.5) {拡散：多い所から少ない所へ};
  \end{scope}
\end{tikzpicture}
\bkfigcaption{2つの流れ。左：電場に押されて、電子がそろって動く（ドリフト）。右：電場がなくても、電子の多い左から少ない右へ、全体として広がる（拡散）。}\label{ph:原-9:fig-drift}
\end{bkfigure}

\phboxhead{アインシュタインの関係}
移動度と拡散係数は、どちらも「電子が結晶の中でぶつかりながら進む」ことから決まるので、別々の量ではありません。2つの間には
\pheq{D=\mu\,\frac{k_{\mathrm B}T}{e}}{cm$^2$/s ＝ cm$^2$/(V·s) × V}
\label{ph:原-9:einstein}の関係があります（\textbf{アインシュタインの関係}。⑪で導きます）。\bkidx{あいんしゅたいんのかんけい}{アインシュタインの関係}300 K で $k_{\mathrm B}T/e=\qty{0.0259}{V}$ なので、ケイ素の電子は $D_{\mathrm n}=1350\times0.0259\approx\qty{35}{cm^2/s}$ です。

\phboxhead{少数キャリアは、拡散距離ほど進んで消える}
p形の中に電子（少数キャリア）を入れると、電子は拡散しながら、周りのたくさんの正孔と出会って消えていきます。消えるまでの平均の時間を\textbf{寿命} $\tau$、その間に拡散で進む平均の距離を\textbf{拡散距離} $L$ といい、
\pheq{L=\sqrt{D\tau}}{cm ＝ $\sqrt{\text{cm}^2\text{/s}\times\text{s}}$}
\label{ph:原-9:L}です。ケイ素の電子で $\tau=\qty{1}{\micro s}$ なら、$L\approx\qty{59}{\micro m}$（例題6）。距離が時間の平方根に比例するのは、拡散が「酔っぱらいの千鳥足」のような、でたらめな歩みの積み重ねだからです（数学 8-13 の熱方程式と同じ形）。

\begin{phbaisuru}[4倍]{寿命 $\tau$}
寿命を4倍にすると、拡散距離は $\sqrt4=2$ 倍にしかなりません。2倍遠くまで行くには、4倍の時間がかかるのです。太陽電池では、光で生まれた電子が消える前に接合（⑥）までたどり着く必要があるので、結晶の欠陥を減らして寿命を長くし、拡散距離を板の厚さ（\qty{0.2}{mm} ほど）より長くするくふうをしています。
\end{phbaisuru}

\Needspace{8\baselineskip}\bksubsection{6}{pn 接合：空乏層と内蔵電位}
\label{ph:原-9:sec6}
1つのケイ素の結晶の中で、左半分を p形、右半分を n形にしたものを\textbf{pn 接合}といいます。\bkidx{ぴーえぬせつごう}{pn 接合}2つの板をはり合わせるのではなく、1枚の結晶に、場所によって違う原子を混ぜて作ります。つないだ瞬間に何が起きるかを、図\ref{ph:原-9:fig-pn}で順に見ます。

\begin{enumerate}[label={\bkcircled{\arabic*}}, leftmargin=2.4\zw, itemsep=1pt, topsep=2pt]
\item n形には電子が多く、p形には正孔が多い。境目をはさんで数が違うので、電子は右から左へ、正孔は左から右へ拡散する。
\item 境目の近くで、電子と正孔が出会って消える。すると、境目の近くの n形には、電子を失ったドナー（＋のイオン）だけが、p形には、電子を受けとったアクセプター（－のイオン）だけが残る。イオンは結晶に固定されていて動けない。
\item 境目の両側に、＋と－の電荷の層ができ、その間に、n形から p形へ向かう電場 $\mathcal E$ が生まれる。この電場は、電子を右へ、正孔を左へ押しもどす向き。
\item 拡散で出ていこうとする流れと、電場で押しもどす流れがつり合った所で、移動が止まる。
\end{enumerate}
キャリアがほとんどいなくなった境目の層を\textbf{空乏層}といいます。\bkidx{くうぼうそう}{空乏層}空乏層の両側には電位の差ができ、これを\textbf{内蔵電位} $V_{\mathrm{bi}}$（bi は built-in＝作りつけの）といいます。\bkidx{ないぞうでんい}{内蔵電位}

\begin{bkfigure}
\begin{tikzpicture}[x=1mm, y=1mm]
  % 上：結晶と電荷
  \fill[bkRed!8] (0,30) rectangle (60,50);
  \fill[bkBlue!8] (60,30) rectangle (120,50);
  \fill[bkGray!25] (48,30) rectangle (72,50);
  \draw[bkInk, line width=0.6pt] (0,30) rectangle (120,50);
  \draw[bkInk, line width=0.5pt, dashed] (60,30) -- (60,50);
  \foreach \x in {6,14,22,30,38} \foreach \y in {35,44} {\txjho{(\x,\y)}}
  \foreach \x in {82,90,98,106,114} \foreach \y in {35,44} {\txjel{(\x,\y)}}
  \foreach \x in {51,56} \foreach \y in {35,40,45} {\draw[bkRed!70!black, line width=0.6pt] (\x,\y) circle (1.5); \node[font=\bfseries\scriptsize, text=bkRed!70!black, inner sep=0pt] at (\x,\y) {$-$};}
  \foreach \x in {64,69} \foreach \y in {35,40,45} {\draw[bkBlue!70!black, line width=0.6pt] (\x,\y) circle (1.5); \node[font=\bfseries\scriptsize, text=bkBlue!70!black, inner sep=0pt] at (\x,\y) {$+$};}
  \node[txj lab, anchor=south] at (24,50.8) {p形};
  \node[txj lab, anchor=south] at (96,50.8) {n形};
  \node[txj lab, anchor=south] at (60,50.8) {空乏層};
  \draw[bkGreen!50!black, line width=1.2pt, -{Stealth[length=2.2mm]}] (70,27) -- (50,27);
  \node[txj lab, text=bkGreen!40!black, anchor=north] at (60,25.5) {電場 $\mathcal E$};
  \draw[txj dim] (48,55.5) -- (72,55.5);
  \node[txj lab, anchor=south] at (60,56) {$W$};
  % 下：電子のエネルギーの帯
  \draw[txj band] (0,14) -- (46,14) .. controls (56,14) and (64,4) .. (74,4) -- (120,4);
  \draw[txj band] (0,0) -- (46,0) .. controls (56,0) and (64,-10) .. (74,-10) -- (120,-10);
  \draw[bkPurple!80!black, line width=0.9pt, dashed] (0,1.8) -- (120,1.8);
  \node[txj lab, text=bkPurple!80!black, anchor=south] at (24,3) {フェルミ準位 $E_{\mathrm F}$（破線）};
  \node[txj lab, anchor=south west] at (1,14.5) {伝導帯の下のはし $E_{\mathrm c}$};
  \node[txj lab, anchor=north east] at (120,-11) {価電子帯の上のはし $E_{\mathrm v}$};
  \foreach \x in {82,92,102,112} {\txjel{(\x,6.4)}}
  \foreach \x in {8,18,28,38} {\txjho{(\x,-2.4)}}
  \draw[txj dim] (124,4) -- (124,14);
  \node[txj lab, anchor=west] at (125,9) {$eV_{\mathrm{bi}}$};
  \draw[bkGray, line width=0.4pt, densely dotted] (74,14) -- (126,14);
\end{tikzpicture}
\bkfigcaption{つり合った pn 接合。上：境目の両側に、動けないイオン（丸で囲んだ＋と－）だけの空乏層（灰色）ができ、n形から p形へ向かう電場が生まれる。下：電子のエネルギーの帯。n形の電子から見ると、p形の伝導帯は高さ $eV_{\mathrm{bi}}$ の坂の上にある。フェルミ準位（破線。⑪）は、つり合いでは左右で同じ高さにそろう。}\label{ph:原-9:fig-pn}
\end{bkfigure}

\noindent 帯の図（図\ref{ph:原-9:fig-pn}の下）で見ると、n形の電子にとって、p形の伝導帯は、高さ $eV_{\mathrm{bi}}$ の「坂の上」にあります。正孔は、帯の図では上下を逆にして考える（正孔は上のほうが低いエネルギー）ので、p形の正孔にとっても、n形は坂の上です。拡散でのぼってこようとする電子と、坂で押しもどされる電子がつり合っているのです。⑪の計算から、
\pheq{V_{\mathrm{bi}}=\frac{k_{\mathrm B}T}{e}\ln\frac{N_{\mathrm a}N_{\mathrm d}}{n_{\mathrm i}^{\,2}}}{V ＝ V × 1}
\label{ph:原-9:vbi}となります。$N_{\mathrm a}$ は p形のアクセプター、$N_{\mathrm d}$ は n形のドナーの数密度です。

\begin{phyomi}[$V_{\mathrm{bi}}=(k_{\mathrm B}T/e)\ln(N_{\mathrm a}N_{\mathrm d}/n_{\mathrm i}^{\,2})$]
「内蔵電位は、坂の両側の電子の数の比の自然対数（ln）× 熱の電圧」と読みます。n形の電子の数は $N_{\mathrm d}$、p形の電子の数は $n_{\mathrm i}^{\,2}/N_{\mathrm a}$ なので、比は $N_{\mathrm a}N_{\mathrm d}/n_{\mathrm i}^{\,2}$ です。この比が $e$ 倍になるごとに、坂が $k_{\mathrm B}T/e=\qty{0.0259}{V}$ ずつ高くなります。どちらも $10^{16}$ cm$^{-3}$ なら、比は $10^{12}$ で、$V_{\mathrm{bi}}\approx\qty{0.72}{V}$ です（例題7）。不純物を10倍に増やしても、比は100倍で、$V_{\mathrm{bi}}$ は $0.0259\times\ln100\approx\qty{0.12}{V}$ 増えるだけです。
\end{phyomi}

\begin{phmijika}[内蔵電位で、電池にはならない]
内蔵電位は 0.7 V もあるのに、pn 接合を導線でつないでも、電流は流れません。導線と半導体のつなぎ目にも電位の差ができて、回路を1周すると、ちょうど打ち消し合うからです。つり合っている物から、エネルギーはとり出せません。光を当てたとき（⑩）だけ、つり合いがくずれて、電流がとり出せます。
\end{phmijika}

\noindent 空乏層の厚さ $W$ は、ケイ素で \qty{1}{\micro m} ほど（髪の毛の太さの約100分の1）です。この中に 0.7 V がかかるので、電場は $\qty{1e4}{V/cm}$（$\qty{1e6}{V/m}$）を超えます（例題7。式は⑫）。空気なら放電が起きるほどの強い電場が、ふつうの部品の中にあるのです。

\Needspace{8\baselineskip}\bksubsection{7}{整流：ダイオードは一方通行の弁}
\label{ph:原-9:sec7}
pn 接合に電圧 $V$ をかけます。p形の側を＋にする向きを\textbf{順方向}、反対を\textbf{逆方向}といいます。\bkidx{じゅんほうこう}{順方向}

\begin{itemize}[label={\color{bkTheme}●}, leftmargin=1.6\zw, itemsep=1pt, topsep=2pt]
\item \textbf{順方向}：外からの電圧が内蔵電位を打ち消す向きなので、坂が $eV$ だけ低くなる（図\ref{ph:原-9:fig-bias}の左）。坂をこえられる電子の数は、ボルツマン因子 $e^{-(\text{坂の高さ})/(k_{\mathrm B}T)}$ に比例するので、坂を $eV$ 下げると、$e^{eV/(k_{\mathrm B}T)}$ 倍に増える。こえた電子は p形に入り、正孔と出会って消えながら進む。これが電流になる。
\item \textbf{逆方向}：坂が $eV$ だけ高くなり、n形の電子も p形の正孔も、坂をこえられない。流れるのは、p形の中にわずかにいる電子（少数キャリア）が坂を転がり落ちる分だけで、電圧によらず、とても小さい。
\end{itemize}

\begin{bkfigure}
\begin{tikzpicture}[x=1mm, y=1mm]
  \foreach \ox/\dh/\ttl/\yl in {0/-4/{順方向（坂が低くなる）}/4.5, 68/5/{逆方向（坂が高くなる）}/8.5} {
    \begin{scope}[shift={(\ox,0)}]
      \draw[txj band] (0,{10+\dh}) -- (20,{10+\dh}) .. controls (26,{10+\dh}) and (32,0) .. (38,0) -- (58,0);
      \draw[bkGray, line width=0.6pt, densely dotted] (0,10) -- (20,10) .. controls (26,10) and (32,0) .. (38,0);
      \node[txj lab, anchor=north] at (29,-7) {\ttl};
      \node[txj lab, anchor=north] at (8,\yl) {p形};
      \node[txj lab, anchor=north] at (50,-1.5) {n形};
    \end{scope}}
  % 順方向：電子がこえる
  \foreach \x in {40,44,48,52,56} {\txjel{(\x,2.3)}}
  \txjel{(30,7)} \txjel{(16,8.3)}
  \draw[bkBlue!70!black, line width=0.9pt, -{Stealth[length=2mm]}] (40,5) .. controls (34,9) and (26,9) .. (18,8.5);
  \node[txj lab, text=bkBlue!70!black, anchor=south] at (30,10.5) {多くこえる};
  % 逆方向
  \foreach \x in {108,112,116,120,124} {\txjel{(\x,2.3)}}
  \draw[bkBlue!70!black, line width=0.9pt, -{Stealth[length=2mm]}] (106,3) .. controls (103,7) and (100,9) .. (100,11);
  \draw[bkRed!75!black, line width=1.1pt] (99,7) -- (103,11) (99,11) -- (103,7);
  \node[txj lab, text=bkBlue!70!black, anchor=south] at (97,16.5) {こえられない};
\end{tikzpicture}
\bkfigcaption{電圧をかけた pn 接合の伝導帯の下のはし（実線。点線は電圧をかけないとき）。左：順方向では坂が $eV$ 下がり、n形の電子がたくさん p形へこえる。右：逆方向では坂が $eV$ 上がり、電子はこえられない。}\label{ph:原-9:fig-bias}
\end{bkfigure}

\noindent 2つを1つの式にまとめたものが、\textbf{ダイオードの式}（ショックレーの式）です。\bkidx{だいおーどのしき}{ダイオードの式}
\pheq{I=I_{\mathrm s}\left(e^{eV/(k_{\mathrm B}T)}-1\right)}{A ＝ A × 1}
\label{ph:原-9:diode}$I_{\mathrm s}$ を\textbf{飽和電流}といい、逆方向に流れる小さな電流の大きさです。小さなダイオードで $\qty{1e-14}{A}$ ほどです。式の $-1$ は、逆方向に流れる一定の分を引いたもので、$V=0$ で $I=0$ になるようにしています。

\begin{phyomi}[$I=I_{\mathrm s}(e^{eV/(k_{\mathrm B}T)}-1)$]
「電圧が $k_{\mathrm B}T/e=\qty{25.9}{mV}$ 増えるごとに、電流は $e\approx2.718$ 倍」と読みます。10倍にするには、$\ln10\times25.9\approx\qty{59.6}{mV}$、つまり約 \qty{0.06}{V} 上げればよい。これが数学 3-12 で見た「ダイオードは約 0.06 V ごとに電流10倍」の正体です。逆方向では $e^{eV/(k_{\mathrm B}T)}$ がすぐ0になり、$I\approx-I_{\mathrm s}$ のまま、ほとんど流れません。
\end{phyomi}

\begin{bkfigure}
\begin{tikzpicture}
\begin{axis}[bkplot, axis lines=left, name=lin, width=54mm, height=44mm, xmin=-0.42, xmax=0.78, ymin=-0.5, ymax=11, xtick={-0.4,0,0.4}, xticklabels={$-0.4$,0,0.4}, ytick={0,2,4,6,8,10}, xlabel={$V$（V）}, ylabel={$I$（mA）}]
  \addplot[bkcurve, domain=-0.42:0.718, samples=200] {1e-11*(exp(x/0.0259)-1)};
  \draw[bkGray, line width=0.4pt, densely dotted] (axis cs:-0.42,0) -- (axis cs:0.78,0);
  \node[font=\sffamily\footnotesize, anchor=east, align=right] at (axis cs:0.55,6.5) {0.6 V あたりから\\急に流れる};
  \node[font=\sffamily\footnotesize, anchor=south west] at (axis cs:-0.36,0.8) {逆はほぼ0};
\end{axis}
\begin{axis}[bkplot, axis lines=left, at={(lin.east)}, anchor=west, xshift=30mm, ymode=log, width=50mm, height=44mm, xmin=0.45, xmax=0.8, ymin=1e-3, ymax=1e3, xtick={0.5,0.6,0.7,0.8}, xticklabels={0.5,0.6,0.7,0.8}, ytick={1e-3,1e-1,1e1,1e3}, yticklabels={0.001,0.1,10,1000}, xlabel={$V$（V）}, ylabel={$I$（mA）}]
  \addplot[bkcurve2, domain=0.48:0.79, samples=80] {1e-11*(exp(x/0.0259)-1)};
  \draw[bkInk, line width=0.6pt] (axis cs:0.6,0.115) -- (axis cs:0.6596,0.115) -- (axis cs:0.6596,1.15);
  \node[font=\sffamily\footnotesize, anchor=north west] at (axis cs:0.6,0.1) {60 mV で10倍};
\end{axis}
\end{tikzpicture}
\bkfigcaption{ダイオードの電流と電圧（$I_{\mathrm s}=\qty{1e-14}{A}$、300 K）。左：ふつうのグラフ。0.6 V あたりから急に流れ始め、逆方向にはほとんど流れない。右：縦軸を対数にした片対数グラフ（数学 3-12）。順方向は直線になり、約 60 mV ごとに電流が10倍。}\label{ph:原-9:fig-iv}
\end{bkfigure}

\noindent ふつうのグラフ（図\ref{ph:原-9:fig-iv}の左）では、ある電圧までほとんど流れず、そこから急に流れ始めるように見えます。ケイ素のダイオードで、この「立ち上がり」はおよそ 0.6〜0.7 V です。立ち上がりは、ほんとうに電流が0から急に変わるのではなく、指数関数がふつうの目もりではそう見えるだけです（右の片対数グラフでは、なめらかな直線）。

\begin{phbaisuru}[10倍]{ダイオードに流す電流}
電流を10倍にしても、電圧は約 \qty{0.06}{V} しか増えません（例題8）。だから、電流を 1 mA から 100 mA まで変えても、ダイオードの電圧は 0.66 V から 0.78 V ほどまでしか変わりません。電子回路の設計では、「ケイ素のダイオードは、流れていれば約 0.7 V」と覚えておいて、まず計算します。反対に、電圧を少し上げすぎると、電流はけた違いに増えるので、LED には電流を決める抵抗を直列に入れるのです（やってみよう）。
\end{phbaisuru}

\phboxhead{整流：交流を一方の向きにそろえる}
ダイオードを交流の電源と負荷の間に入れると、電流は順方向の半分の時間だけ流れ、向きが一方にそろいます。これを\textbf{整流}といいます。\bkidx{せいりゅう}{整流}スマホの充電器（USB の電源）の中では、ダイオード4個をひし形に組んで（ブリッジ）、交流の両方の向きを、同じ向きの電流にしています。物理シリーズの『電磁気学』（電-16 ⑤）で「一方の向きにだけ電流を通す弁」とした部品の中身が、この pn 接合です。

\begin{phmodel}[理想ダイオードの式]
\phignore{半導体の中や電極の抵抗（直列抵抗）。空乏層の中で電子と正孔が出会って消えること。大きな電流で、入ってくるキャリアが多数キャリアと同じくらいになること。}
\phassume{接合は段の形（境目で p形から n形へ急に変わる）。空乏層の外には電場がない。坂をこえたキャリアは、拡散だけで進む。}
\phlimit{大きな電流（数十 mA から上）では直列抵抗で電圧が上がる。小さな電流では、空乏層の中での再結合で、10倍の電圧が 60 mV より大きくなる（120 mV に近づく）。逆方向に大きな電圧をかけると、急に電流が流れる（降伏）。}
\end{phmodel}

\Needspace{14\baselineskip}\bksubsection{8}{光を出す・光を吸う：直接遷移と間接遷移}
\label{ph:原-9:sec8}
伝導帯の電子が、価電子帯の正孔に落ちて再結合するとき、電子はエネルギー $E_{\mathrm g}$ くらいを手放します。手放し方は2通りです。
\begin{itemize}[label={\color{bkTheme}●}, leftmargin=1.6\zw, itemsep=1pt, topsep=2pt]
\item 光子を1個出す。光子のエネルギーは $hf\approx E_{\mathrm g}$（原-4 の「準位の差が光になる」と同じ）。
\item 結晶の原子のゆれ（熱）として手放す。
\end{itemize}
どちらになりやすいかは、材料で決まります。原-8 ⑦で見たように、結晶の中の電子は波で、エネルギーのほかに「波数」をもっています。ガリウムヒ素（GaAs）や窒化ガリウム（GaN）では、伝導帯のいちばん下と価電子帯のいちばん上が、同じ波数の所にあるので、電子はまっすぐ落ちて光子を出せます（\textbf{直接遷移}）。ケイ素では、2つが別の波数の所にあり、落ちるには原子のゆれに波数を受けとってもらう必要があります（\textbf{間接遷移}）。これは起こりにくいので、その前に、結晶の欠陥などを通って熱に変わってしまいます。\bkidx{ちょくせつせんい}{直接遷移}

\begin{bkfigure}
\begin{tikzpicture}[x=1mm, y=1mm]
  \begin{scope}
    \draw[bkGray, line width=0.5pt, -{Stealth[length=1.6mm]}] (0,0) -- (48,0);
    \node[txj lab, anchor=north] at (46,-0.8) {波数};
    \draw[bkGray, line width=0.5pt, -{Stealth[length=1.6mm]}] (24,-2) -- (24,40);
    \node[txj lab, anchor=south] at (24,40.5) {エネルギー};
    \draw[txj band, domain=-16:16, samples=40] plot ({24+\x}, {22+0.05*\x*\x});
    \draw[txj band, domain=-16:16, samples=40] plot ({24+\x}, {12-0.05*\x*\x});
    \txjel{(24,23.4)} \txjho{(24,10.6)}
    \draw[bkOrange!85!black, line width=1.1pt, -{Stealth[length=2mm]}] (24,21.5) -- (24,12.8);
    \draw[bkOrange!85!black, line width=0.8pt, decorate, decoration={snake, amplitude=0.6mm, segment length=2.2mm}] (26,17) -- (38,17);
    \draw[bkOrange!85!black, line width=0.8pt, -{Stealth[length=1.6mm]}] (38,17) -- (41,17);
    \node[txj lab, text=bkOrange!75!black, anchor=south] at (36,18) {光子};
    \node[txj lab, anchor=north] at (24,-5) {直接遷移（GaAs・GaN）};
  \end{scope}
  \begin{scope}[shift={(66,0)}]
    \draw[bkGray, line width=0.5pt, -{Stealth[length=1.6mm]}] (0,0) -- (52,0);
    \node[txj lab, anchor=north] at (50,-0.8) {波数};
    \draw[bkGray, line width=0.5pt, -{Stealth[length=1.6mm]}] (14,-2) -- (14,40);
    \node[txj lab, anchor=south] at (14,40.5) {エネルギー};
    \draw[txj band, domain=-12:12, samples=40] plot ({36+\x}, {22+0.06*\x*\x});
    \draw[txj band, domain=-12:12, samples=40] plot ({14+\x}, {12-0.06*\x*\x});
    \txjel{(36,23.4)} \txjho{(14,10.6)}
    \draw[bkGray, line width=0.9pt, dashed, -{Stealth[length=2mm]}] (34.5,21.5) -- (15.5,12.5);
    \node[txj lab, anchor=west, align=left] at (27,13) {原子のゆれの\\助けが要る};
    \node[txj lab, anchor=north] at (26,-5) {間接遷移（ケイ素）};
  \end{scope}
\end{tikzpicture}
\bkfigcaption{エネルギーと波数でかいた帯の形（ふちだけ）。左：伝導帯の底と価電子帯の頂上が同じ波数にあり、電子はまっすぐ落ちて光子を出す。右：ケイ素では波数がずれていて、落ちるには原子のゆれの助けが要り、光になりにくい。}\label{ph:原-9:fig-direct}
\end{bkfigure}

\noindent 反対に、光を吸うときも、エネルギーが $E_{\mathrm g}$ より大きい光子だけが、電子を価電子帯から伝導帯へ上げられます。ケイ素の $E_{\mathrm g}=\qty{1.12}{eV}$ に当たる波長は $1240/1.12\approx\qty{1110}{nm}$ で、これより長い波長の赤外線は、ケイ素を素通りします。ケイ素は間接遷移なので、光を吸う力が弱く、太陽電池では \qty{0.1}{mm} 以上の厚さを使います（光を吸うときも、原子のゆれの助けが要るため）。

\Needspace{14\baselineskip}\bksubsection{9}{LED：はてな？の答え}
\label{ph:原-9:sec9}
直接遷移の半導体で pn 接合を作り、順方向に電圧をかけます。n形から電子が、p形から正孔が、接合に流れこみます。接合の所で電子と正孔が出会い、電子が $E_{\mathrm g}$ だけ落ちて、光子を出します（図\ref{ph:原-9:fig-led}）。これが\textbf{発光ダイオード}（LED：light emitting diode）です。\bkidx{はっこうだいおーど}{発光ダイオード}

\begin{bkfigure}
\begin{tikzpicture}[x=1mm, y=1mm]
  \fill[bkRed!8] (0,-12) rectangle (44,20);
  \fill[bkBlue!8] (44,-12) rectangle (88,20);
  \node[txj lab, anchor=south] at (20,20.5) {p形};
  \node[txj lab, anchor=south] at (68,20.5) {n形};
  \draw[txj band] (0,13) -- (36,13) .. controls (40,13) and (48,10) .. (52,10) -- (88,10);
  \draw[txj band] (0,1) -- (36,1) .. controls (40,1) and (48,-2) .. (52,-2) -- (88,-2);
  \foreach \x in {58,66,74,82} {\txjel{(\x,12.3)}}
  \foreach \x in {6,14,22,30} {\txjho{(\x,-1.3)}}
  \draw[bkBlue!70!black, line width=0.9pt, -{Stealth[length=2mm]}] (56,14.5) -- (47,14.5);
  \draw[bkRed!70!black, line width=0.9pt, -{Stealth[length=2mm]}] (32,-3.5) -- (41,-3.5);
  \txjel{(44,11.3)} \txjho{(44,-0.3)}
  \draw[bkOrange!85!black, line width=1.1pt, -{Stealth[length=2mm]}] (44,9.6) -- (44,1.4);
  \draw[bkOrange!85!black, line width=0.8pt, decorate, decoration={snake, amplitude=0.6mm, segment length=2.2mm}] (46,5.5) -- (58,5.5);
  \draw[bkOrange!85!black, line width=0.8pt, -{Stealth[length=1.6mm]}] (58,5.5) -- (61,5.5);
  \node[txj lab, text=bkOrange!75!black, anchor=west] at (62,5.5) {光子 $hf\approx E_{\mathrm g}$};
  \node[txj lab, anchor=north] at (44,-12.8) {順方向に電圧をかけて、電子と正孔を接合に流しこむ};
\end{tikzpicture}
\bkfigcaption{LED の帯の図。順方向で坂が低くなり、n形から電子（左向きの矢印）、p形から正孔（右向きの矢印）が接合に入る。出会うと電子が $E_{\mathrm g}$ だけ落ち、そのエネルギーが1個の光子になる。}\label{ph:原-9:fig-led}
\end{bkfigure}

\noindent はてな？の答えです。LED は、電子1個のエネルギーを、熱をはさまずに、そのまま光子1個に渡しています。白熱電球は、電気でフィラメントを熱し、熱した物が出す光（熱-13）を使うので、エネルギーの大部分が目に見えない赤外線と熱になります。LED は熱を通らないので、同じ明るさを、白熱電球の数分の1の電気で出せるのです。

光の色は、光子のエネルギー $hf\approx E_{\mathrm g}$ で決まります。原-2 の $E=hc/\lambda$ から、
\pheq{\lambda\approx\frac{hc}{E_{\mathrm g}}\qquad\Bigl(\lambda\,\text{〔nm〕}\approx\frac{1240}{E_{\mathrm g}\,\text{〔eV〕}}\Bigr)}{m ＝ J·s × m/s ÷ J}
\label{ph:原-9:color}です。

\begin{bkcenter}
\begin{bktable}{colspec={Q[l,m,wd=32mm] Q[c,m,wd=18mm] Q[c,m,wd=20mm] Q[l,m,wd=48mm]}, rowsep=1pt, row{1}={bg=bkTheme!10}, cells={font=\small}}
材料（主なもの） & $E_{\mathrm g}$〔eV〕 & 波長〔nm〕 & 色・使い道 \\
InGaN（窒化物） & 約2.7 & 約460 & 青（白色の LED のもと） \\
InGaN（窒化物） & 約2.3 & 約540 & 緑（信号機の「青」） \\
AlGaInP・GaAsP & 約1.9 & 約650 & 赤 \\
GaAs & 1.42 & 約870 & 赤外（リモコン・センサー） \\
ケイ素（間接遷移） & 1.12 & 約1110 & ほとんど光らない \\
\end{bktable}
\bktabcaption{LED の材料と色の目安（波長は $1240/E_{\mathrm g}$。実際の LED の光は、温度などで少しずれ、はばをもつ）}\label{ph:原-9:tab-color}
\end{bkcenter}

\begin{phbaisuru}{禁制帯の幅 $E_{\mathrm g}$}
禁制帯の幅を大きくすると、光子のエネルギーが大きくなり、波長は短く（赤→緑→青→紫外）なります。$E_{\mathrm g}$ が2倍なら、波長は半分です。光らせるには、少なくとも $E_{\mathrm g}/e$ くらいの電圧が要るので、青い LED（約 2.7〜\qty{3}{V}）は赤（約 1.9〜\qty{2}{V}）より大きな電圧が要ります。やってみようで、同じ 3 V の電池と \qty{100}{\ohm} で青が暗かったのは、抵抗にかかる電圧が $3-1.9=\qty{1.1}{V}$ から 0.2 V ほどに減り、電流が約 \qty{11}{mA} から約 \qty{2}{mA} に減ったからです。
\end{phbaisuru}

\phboxhead{青い LED と白い光}
赤と緑の LED は 1960〜70年代に作られましたが、青は、大きな禁制帯をもつ窒化ガリウムのよい結晶と、p形がなかなか作れず、長いあいだ実現しませんでした。赤崎・天野・中村の3人は、1980〜90年代に、窒化ガリウムの結晶の育て方と p形の作り方を見つけ、明るい青色 LED を実現しました（2014年ノーベル物理学賞「明るく省エネルギーな白色光源を可能にした、効率のよい青色発光ダイオードの発明」）。今の白い LED の多くは、青い LED の光の一部を黄色に光る物質（蛍光体）に当て、青と黄色を混ぜて白に見せています。

\Needspace{16\baselineskip}\bksubsection{10}{太陽電池とフォトダイオード}
\label{ph:原-9:sec10}
LED と逆向きのことをするのが\textbf{太陽電池}です。\bkidx{たいようでんち}{太陽電池}pn 接合に、$E_{\mathrm g}$ より大きなエネルギーの光子が当たると、電子が価電子帯から伝導帯へ上がり、電子と正孔が1組できます（図\ref{ph:原-9:fig-solar}）。空乏層の電場（⑥）は、電子を n形の側へ、正孔を p形の側へ押し分けます。空乏層の外で生まれた組も、拡散距離（⑤）の内側なら、拡散で空乏層にたどり着き、押し分けられます。n形には電子が、p形には正孔がたまり、n形が－、p形が＋の電池になります。これが、光電池がモーターを回すしくみです。

\noindent 外に流れる電流は、光で押し分けられる分 $I_{\mathrm L}$（\textbf{光電流}）から、電池の電圧 $V$ のせいで pn 接合を順方向に流れて戻ってしまう分（⑦のダイオードの電流）を引いたものです。
\pheq{I=I_{\mathrm L}-I_{\mathrm s}\left(e^{eV/(k_{\mathrm B}T)}-1\right)}{A ＝ A － A}
\label{ph:原-9:solar}光電流 $I_{\mathrm L}$ は、1秒に吸われて押し分けられる光子の数に $e$ をかけたものなので、光の強さに比例します（①の「光を2倍で電流2倍」）。

\begin{bkfigure}
\begin{tikzpicture}[x=1mm, y=1mm]
  \begin{scope}
    \fill[bkRed!8] (0,-6) rectangle (30,22);
    \fill[bkBlue!8] (30,-6) rectangle (60,22);
    \node[txj lab, anchor=south] at (8,22.5) {p形};
    \node[txj lab, anchor=south] at (48,22.5) {n形};
    \draw[txj band] (0,16) -- (22,16) .. controls (27,16) and (33,8) .. (38,8) -- (60,8);
    \draw[txj band] (0,4) -- (22,4) .. controls (27,4) and (33,-4) .. (38,-4) -- (60,-4);
    \draw[bkOrange!85!black, line width=0.8pt, decorate, decoration={snake, amplitude=0.6mm, segment length=2.2mm}] (22,32) -- (27,15);
    \draw[bkOrange!85!black, line width=0.8pt, -{Stealth[length=1.6mm]}] (27,15) -- (27.6,13);
    \node[txj lab, text=bkOrange!75!black, anchor=east] at (21,31) {光子};
    \txjel{(30,13)} \txjho{(30,1)}
    \draw[bkGray, line width=0.6pt, dashed] (30,11.6) -- (30,2.4);
    \draw[bkBlue!70!black, line width=0.9pt, -{Stealth[length=2mm]}] (32,12.2) -- (40,9.6);
    \draw[bkRed!70!black, line width=0.9pt, -{Stealth[length=2mm]}] (28,1.6) -- (20,4.6);
    \node[txj lab, anchor=north] at (30,-7) {光で生まれた電子と正孔が、坂で分かれる};
  \end{scope}
  \begin{scope}[shift={(74,2)}]
    \draw[bkInk, line width=0.7pt] (0,0) -- (0,18) -- (44,18) -- (44,0) -- cycle;
    \fill[white] (-4,5) rectangle (4,13);
    \draw[bkInk, line width=0.7pt, fill=bkRed!8] (-4,5) rectangle (4,9);
    \draw[bkInk, line width=0.7pt, fill=bkBlue!8] (-4,9) rectangle (4,13);
    \node[font=\sffamily\footnotesize] at (0,7) {p};
    \node[font=\sffamily\footnotesize] at (0,11) {n};
    \txjR{(22,18)}{0}
    \node[txj lab, anchor=south] at (22,20.5) {負荷};
    \draw[bkInk, txj arr] (14,0) -- (30,0);
    \node[txj lab, anchor=north] at (22,-1.5) {電流 $I$};
    \draw[bkOrange!85!black, line width=0.8pt, decorate, decoration={snake, amplitude=0.6mm, segment length=2.2mm}] (-16,9) -- (-6,9);
    \draw[bkOrange!85!black, line width=0.8pt, -{Stealth[length=1.6mm]}] (-6,9) -- (-4.6,9);
    \node[txj lab, anchor=north] at (0,-1.5) {セル};
  \end{scope}
\end{tikzpicture}
\bkfigcaption{太陽電池。左：空乏層で生まれた電子は n形へ、正孔は p形へ、坂で分けられる。右：p形が＋、n形が－の電池になり、外の負荷に、p形から出る向きの電流が流れる（ダイオードとしては逆方向の向き）。}\label{ph:原-9:fig-solar}
\end{bkfigure}

\begin{phyomi}[$I=I_{\mathrm L}-I_{\mathrm s}(e^{eV/(k_{\mathrm B}T)}-1)$]
「光が押し出す電流から、接合をもどる電流を引いたものが、外に出る電流」と読みます。導線で短くつなぐ（$V=0$）と、$I=I_{\mathrm L}$（\textbf{短絡電流}）。何もつながない（$I=0$）と、もどる電流が $I_{\mathrm L}$ と等しくなる電圧
\[ V_{\mathrm{oc}}=\frac{k_{\mathrm B}T}{e}\ln\Bigl(\frac{I_{\mathrm L}}{I_{\mathrm s}}+1\Bigr) \]
になります（\textbf{開放電圧}）。\bkidx{かいほうでんあつ}{開放電圧}ケイ素の太陽電池1枚（セル）で、0.6〜0.7 V ほどです。
\end{phyomi}

\phboxhead{いちばん大きな電力と、曲線因子・効率}
太陽電池から取り出せる電力は $P=VI$ です。短絡（$V=0$）でも開放（$I=0$）でも、電力は0です。その間に、$P$ がいちばん大きくなる電圧があります（図\ref{ph:原-9:fig-iv-solar}）。数学 5-7 のように $\dd P/\dd V=0$ で求めてもよいし、表計算で電圧を少しずつ変えて、$VI$ のいちばん大きな所を探してもかまいません（例題10）。いちばん大きな電力 $P_{\max}$ と、$V_{\mathrm{oc}}\times I_{\mathrm L}$（グラフの長方形の面積）との比を\textbf{曲線因子}（FF）といいます。\bkidx{きょくせんいんし}{曲線因子}電流と電圧のグラフが、どれだけ角ばっているかを表し、よいセルで 0.8 ほどです。太陽電池の\textbf{効率}は、当たった光の仕事率に対する $P_{\max}$ の割合です。性能を比べるときは、1 m$^2$ に \qty{1000}{W} の光を当ててはかります（晴れた日の昼、地上で受ける太陽の光の強さの目安。大気の外では 1 m$^2$ に \qty{1361}{W}：原-2）。

\begin{bkfigure}
\begin{tikzpicture}
\begin{axis}[bkplot, axis lines=left, width=100mm, height=56mm, xmin=0, xmax=0.7, ymin=0, ymax=7.6, xtick={0,0.1,0.2,0.3,0.4,0.5,0.6,0.7}, xticklabels={0,0.1,0.2,0.3,0.4,0.5,0.6,0.7}, ytick={0,1,2,3,4,5,6,7}, xlabel={電圧 $V$（V）}, ylabel={電流 $I$（A）・電力 $P$（W）}]
  \addplot[bkcurve, domain=0:0.6288, samples=150] {3.5-1e-10*(exp(x/0.0259)-1)};
  \addplot[bkcurve, dashed, domain=0:0.6468, samples=150] {7.0-1e-10*(exp(x/0.0259)-1)};
  \addplot[bkcurve2, domain=0:0.6288, samples=150] {x*(3.5-1e-10*(exp(x/0.0259)-1))};
  \draw[bkGray, line width=0.6pt, densely dotted] (axis cs:0,3.333) -- (axis cs:0.5485,3.333) -- (axis cs:0.5485,0);
  \fill[bkInk] (axis cs:0.5485,1.833) circle (1.6pt);
  \node[font=\sffamily\footnotesize, anchor=south east] at (axis cs:0.54,1.95) {$P_{\max}=\qty{1.83}{W}$};
  \node[font=\sffamily\footnotesize, text=bkBlue!80!black, anchor=south west] at (axis cs:0.02,3.55) {電流（実線）};
  \node[font=\sffamily\footnotesize, text=bkBlue!80!black, anchor=south west] at (axis cs:0.02,7.05) {光を2倍にした電流（破線）};
  \node[font=\sffamily\footnotesize, text=bkOrange!80!black, anchor=south west] at (axis cs:0.03,1.55) {電力 $P=VI$（山の形の線）};
\end{axis}
\end{tikzpicture}
\bkfigcaption{太陽電池（面積 \qty{100}{cm^2}、例題10）の電流と電圧（実線）と電力（山の形の線）。点線の長方形の面積が $P_{\max}$。光を2倍にすると（破線）、電流は2倍になるが、開放電圧は \qty{18}{mV} しか増えない。}\label{ph:原-9:fig-iv-solar}
\end{bkfigure}

\begin{phbaisuru}{光の強さ}
光を2倍にすると、短絡電流 $I_{\mathrm L}$ は2倍になりますが、開放電圧は $\frac{k_{\mathrm B}T}{e}\ln2\approx\qty{18}{mV}$ しか増えません（$\ln$ の中が2倍になるだけ）。だから、取り出せる電力は、ほぼ光の強さに比例します（例題10 のセルで 1.83 W → 3.79 W）。曇りの日に発電が減るのは、主に電流が減るからです。
\end{phbaisuru}

\phboxhead{効率の上限：ショックレーとクワイサー}
太陽の光には、いろいろな波長がふくまれています。$E_{\mathrm g}$ より小さなエネルギーの光子は吸われず、素通りします。$E_{\mathrm g}$ より大きな光子も、余ったエネルギーは、すぐ熱になってしまい、電子1個から取り出せるのは $E_{\mathrm g}$ より少し小さなエネルギーだけです。この2つの損失を考えると、1つの pn 接合の太陽電池の効率には上限があり、ケイ素の $E_{\mathrm g}$ では約30 \% になります（ショックレーとクワイサー、1961年）。市販の太陽電池パネルの効率は、20 \% 前後から 25 \% ほどです。禁制帯の違う材料を何段も重ねると、上限をこえられ、人工衛星などで使われています。

\phboxhead{フォトダイオード：光をはかる}
同じ pn 接合を、光の強さをはかるのに使うものが\textbf{フォトダイオード}です。\bkidx{ふぉとだいおーど}{フォトダイオード}逆方向に電圧をかけておくと、流れる電流はほぼ光電流 $I_{\mathrm L}$ だけになり、光の強さにきれいに比例します。1 W の光あたりの電流（\textbf{感度}〔A/W〕）は、当たった光子のうち電子を押し出す割合（量子効率 $\eta$）を使って、
\pheq{\text{感度}=\eta\,\frac{e}{hf}=\eta\,\frac{e\lambda}{hc}}{A/W ＝ C ÷ J}
\label{ph:原-9:resp}です（例題12）。デジタルカメラやスマホのカメラの画素、光ファイバー通信の受け手、テレビのリモコンの受光部も、フォトダイオードです。放射線をはかる半導体の検出器（原-14 ②）も、同じしくみです。

\begin{phmijika}[LED は光センサーにもなる]
LED の足をテスターの電圧のはんいにつなぎ、明るい光を当てると、小さな電圧が出ます。LED も pn 接合なので、光を当てれば太陽電池のようにはたらくのです（ただし、自分が出す色より波長の短い光でないと、電子を上げられません）。テレビのリモコンの赤外 LED は、スマホのカメラで見ると光っているのがわかります。カメラのフォトダイオードは、目に見えない赤外線にも反応するからです。
\end{phmijika}

\phkokomade

\Needspace{9\baselineskip}\phstage{shiki}{フェルミ分布から導く}{5}

\bksubsection{11}{キャリアの数の式：状態の数 × フェルミ分布の裾}
\label{ph:原-9:sec11}
③と④で結果だけ使った式を、原-8 の道具で導きます。伝導帯の電子の数は、「エネルギー $E$ の所にある席の数」×「その席が電子でうまっている割合」を、伝導帯の下のはし $E_{\mathrm c}$ から上まで足し合わせたものです。

\phboxhead{席の数：3次元の箱の状態密度}
原-8 ④では、箱の中の自由な電子の状態を数え、エネルギー $E$ から $E+\dd E$ の間にある1 cm$^3$ あたりの状態の数（\textbf{状態密度}）が $\sqrt{E}$ に比例することを見ました。結晶の中の伝導帯の電子も、下のはしの近くでは、質量を「有効質量」$m_{\mathrm n}^{*}$ に置きかえた自由な電子のようにふるまいます。だから、
\[ g_{\mathrm c}(E)=\frac{1}{2\pi^2}\Bigl(\frac{2m_{\mathrm n}^{*}}{\hbar^2}\Bigr)^{3/2}\sqrt{E-E_{\mathrm c}} \]
です（スピンの2倍をふくむ）。

\phboxhead{うまっている割合：フェルミ分布の裾}
うまっている割合は、熱-15 ⑨のフェルミ分布です。熱-15 ではエネルギーを $\varepsilon$、化学ポテンシャルを $\mu$ と書きましたが、固体の本ではそれぞれ $E$、$E_{\mathrm F}$（\textbf{フェルミ準位}）と書きます（原-8 ⑤）。
\[ f(E)=\frac{1}{e^{(E-E_{\mathrm F})/(k_{\mathrm B}T)}+1} \]
半導体では、$E_{\mathrm F}$ は禁制帯の中にあり、伝導帯の席は $E_{\mathrm F}$ より何 $k_{\mathrm B}T$ も上にあります。そこでは分母の指数がとても大きいので、$+1$ を無視して $f\approx e^{-(E-E_{\mathrm F})/(k_{\mathrm B}T)}$（ボルツマン因子）としてよくなります。

\phboxhead{足し合わせる}
$x=(E-E_{\mathrm c})/(k_{\mathrm B}T)$ とおくと、
\begin{align*}
n&=\int_{E_{\mathrm c}}^{\infty}g_{\mathrm c}(E)\,e^{-(E-E_{\mathrm F})/(k_{\mathrm B}T)}\,\dd E\\
&=\frac{1}{2\pi^2}\Bigl(\frac{2m_{\mathrm n}^{*}k_{\mathrm B}T}{\hbar^2}\Bigr)^{3/2}e^{-(E_{\mathrm c}-E_{\mathrm F})/(k_{\mathrm B}T)}\int_0^{\infty}\sqrt{x}\,e^{-x}\,\dd x
\end{align*}
最後の積分は $\sqrt{\pi}/2\approx0.886$ です（$x=u^2$ とおくと、数学 8-6⑦のガウス積分 $\int_{-\infty}^{\infty}e^{-u^2}\dd u=\sqrt\pi$ の半分の形になります）。まとめると、
\pheq{n=N_{\mathrm c}\,e^{-(E_{\mathrm c}-E_{\mathrm F})/(k_{\mathrm B}T)},\qquad N_{\mathrm c}=2\Bigl(\frac{2\pi m_{\mathrm n}^{*}k_{\mathrm B}T}{h^2}\Bigr)^{3/2}}{cm$^{-3}$}
\label{ph:原-9:nc}となります。$N_{\mathrm c}$ は、「伝導帯の下のはしから $k_{\mathrm B}T$ くらいの範囲にある席の数」にあたり、③の $N_{\mathrm c}$ の正体です（$T^{3/2}$ に比例）。正孔も同じように、価電子帯の上のはしから下へ足して、
\pheq{p=N_{\mathrm v}\,e^{-(E_{\mathrm F}-E_{\mathrm v})/(k_{\mathrm B}T)}}{cm$^{-3}$}
です。

\phboxhead{かけ算すると、フェルミ準位が消える}
2つをかけると、$E_{\mathrm c}-E_{\mathrm v}=E_{\mathrm g}$ なので、
\[ np=N_{\mathrm c}N_{\mathrm v}\,e^{-(E_{\mathrm c}-E_{\mathrm v})/(k_{\mathrm B}T)}=N_{\mathrm c}N_{\mathrm v}\,e^{-E_{\mathrm g}/(k_{\mathrm B}T)} \]
となり、$E_{\mathrm F}$ が消えます。まじりけを入れて $E_{\mathrm F}$ が動いても、かけ算は変わりません。これが④の $np=n_{\mathrm i}^{\,2}$ で、平方根をとると③の式 $n_{\mathrm i}=\sqrt{N_{\mathrm c}N_{\mathrm v}}\,e^{-E_{\mathrm g}/(2k_{\mathrm B}T)}$ になります。

\begin{phflow}[フェルミ分布から、キャリアの式へ]
\phflowrow{$n=\displaystyle\int g_{\mathrm c}(E)f(E)\,\dd E$}{$E-E_{\mathrm F}\gg k_{\mathrm B}T$\\（ボルツマンの裾）}{$n=N_{\mathrm c}e^{-(E_{\mathrm c}-E_{\mathrm F})/(k_{\mathrm B}T)}$}
\phflowrow{$n$ と $p$ の式}{2つをかけて\\$E_{\mathrm F}$ を消す}{$np=n_{\mathrm i}^{\,2}$}
\phflowrow{$n=N_{\mathrm c}e^{-(E_{\mathrm c}-E_{\mathrm F})/(k_{\mathrm B}T)}$}{n形で $n=N_{\mathrm d}$\\（ドナーは全部出す）}{$E_{\mathrm c}-E_{\mathrm F}=k_{\mathrm B}T\ln\dfrac{N_{\mathrm c}}{N_{\mathrm d}}$}
\end{phflow}

\noindent 3段めの式で、n形のフェルミ準位の位置がわかります。$N_{\mathrm d}=10^{16}$ なら、$E_{\mathrm F}$ は伝導帯の下のはしより \qty{0.21}{eV} 下です（例題5）。まじりけのない半導体では、$E_{\mathrm F}$ は禁制帯のほぼまん中にあります。

\phboxhead{アインシュタインの関係を導く}
つり合いでは、場所が変わっても電流は0です。電子の流れは、ドリフトと拡散の和です。電子の電荷は $-e$ なので、電流の密度は
\[ j_{\mathrm n}=en\mu_{\mathrm n}\mathcal E+eD_{\mathrm n}\frac{\dd n}{\dd x}=0 \]
（1つめがドリフト、2つめが拡散。電子は坂を下る向きに広がり、電荷が負なので電流はその逆向き）。一方、電位 $\phi$ の所の電子のエネルギーは $-e\phi$ だけずれるので、ボルツマン因子から $n\propto e^{e\phi/(k_{\mathrm B}T)}$ です。微分すると、$\mathcal E=-\dd\phi/\dd x$ を使って
\[ \frac{\dd n}{\dd x}=\frac{e}{k_{\mathrm B}T}\,n\,\frac{\dd\phi}{\dd x}=-\frac{e}{k_{\mathrm B}T}\,n\,\mathcal E \]
これを上の式に入れると、$en\mathcal E\,(\mu_{\mathrm n}-eD_{\mathrm n}/(k_{\mathrm B}T))=0$ となり、
\[ D_{\mathrm n}=\mu_{\mathrm n}\,\frac{k_{\mathrm B}T}{e} \]
が出ます（⑤のアインシュタインの関係）。「電場で押す流れ」と「数の坂で広がる流れ」がいつも打ち消し合うためには、2つの係数がこの比でなければならないのです。

\phboxhead{内蔵電位を導く}
同じボルツマン因子を、pn 接合の両はしに使います。n形の奥では $n=N_{\mathrm d}$、p形の奥では $n=n_{\mathrm i}^{\,2}/N_{\mathrm a}$ です。電位の差が $V_{\mathrm{bi}}$ なので、
\[ \frac{N_{\mathrm d}}{n_{\mathrm i}^{\,2}/N_{\mathrm a}}=e^{eV_{\mathrm{bi}}/(k_{\mathrm B}T)}\qquad\text{より}\qquad V_{\mathrm{bi}}=\frac{k_{\mathrm B}T}{e}\ln\frac{N_{\mathrm a}N_{\mathrm d}}{n_{\mathrm i}^{\,2}} \]
となり、⑥の式が出ます。

\Needspace{8\baselineskip}\bksubsection{12}{空乏層の電場とダイオードの式}
\label{ph:原-9:sec12}
\phboxhead{ポアソン方程式で、空乏層の厚さを求める}
空乏層の中には、動けないイオンの電荷だけがあると考えます（\textbf{空乏近似}）。境目を $x=0$ とし、p形の側の $-x_{\mathrm p}<x<0$ に電荷密度 $\rho_{\mathrm q}=-eN_{\mathrm a}$、n形の側の $0<x<x_{\mathrm n}$ に $\rho_{\mathrm q}=+eN_{\mathrm d}$ があります。全体の電荷は0なので、
\[ N_{\mathrm a}x_{\mathrm p}=N_{\mathrm d}x_{\mathrm n} \]
です。1次元のポアソン方程式（電-2 ⑩。$\varepsilon=\varepsilon_{\mathrm r}\varepsilon_0$、ケイ素の比誘電率 $\varepsilon_{\mathrm r}=11.7$）
\pheq{\frac{\dd\mathcal E}{\dd x}=\frac{\rho_{\mathrm q}}{\varepsilon},\qquad \mathcal E=-\frac{\dd\phi}{\dd x}}{V/cm$^2$ ＝ C/cm$^3$ ÷ (F/cm)}
\label{ph:原-9:poisson}を、空乏層の外で $\mathcal E=0$ として積分します。電荷が一様なので、電場は直線で変わり、境目でいちばん強くなります（図\ref{ph:原-9:fig-depl}）。
\[ |\mathcal E_{\max}|=\frac{eN_{\mathrm d}x_{\mathrm n}}{\varepsilon}=\frac{eN_{\mathrm a}x_{\mathrm p}}{\varepsilon} \]
電位の差は、電場のグラフの三角形の面積なので、$V_{\mathrm{bi}}=\tfrac12|\mathcal E_{\max}|\,W$（$W=x_{\mathrm p}+x_{\mathrm n}$）。この3つの式から $x_{\mathrm p}$・$x_{\mathrm n}$ を消すと、
\pheq{W=\sqrt{\frac{2\varepsilon V_{\mathrm{bi}}}{e}\Bigl(\frac{1}{N_{\mathrm a}}+\frac{1}{N_{\mathrm d}}\Bigr)},\qquad |\mathcal E_{\max}|=\frac{2V_{\mathrm{bi}}}{W}}{cm}
\label{ph:原-9:W}となります。逆方向に電圧 $V_{\mathrm R}$ をかけると、$V_{\mathrm{bi}}$ が $V_{\mathrm{bi}}+V_{\mathrm R}$ に変わり、空乏層は厚くなります。

\phboxhead{プログラムで、電荷から電場と電位を順に足す}
ポアソン方程式は、電-2 ⑪のように、小さな区間に分けて足していけば、どんな電荷の並び方でも解けます。例題7 の接合で、電荷から電場、電場から電位を足していきます。
\begin{phcode}[Python]
e, eps = 1.60e-19, 11.7 * 8.85e-14     \# C、F/cm
Na = Nd = 1e16                          \# cm\^{}-3
xp = xn = 0.2152e-4                     \# cm（例題7の空乏層の半分）
N = 1000
dx = (xp + xn) / N
E, phi = 0.0, 0.0
for i in range(N):
    x = -xp + (i + 0.5) * dx            \# マスのまん中
    rho = -e * Na if x < 0 else e * Nd  \# 空乏層の電荷密度
    E = E + rho / eps * dx              \# dE/dx = rho/eps
    phi = phi - E * dx                  \# 電位は電場の面積だけ下がる
    if (i + 1) \% 250 == 0:
        print(round((x + dx/2) * 1e4, 3), round(E), round(phi, 3))
\end{phcode}
\begin{phcodeout}[実行結果]
-0.108 -16627 0.09
-0.0 -33253 0.359
0.108 -16627 0.627
0.215 0 0.716
\end{phcodeout}
左から、場所（\si{\micro m}）、電場（V/cm）、p形のはしからはかった電位（V）です。電場は境目で $\qty{-3.3e4}{V/cm}$、n形のはしで0にもどり、電位はちょうど $V_{\mathrm{bi}}=\qty{0.716}{V}$ になります。実際の接合のように、不純物がなだらかに変わる場合は、\texttt{rho} の行を変えるだけで計算できます。

\begin{bkfigure}
\begin{tikzpicture}
\begin{axis}[bkplot, axis lines=left, name=rq, width=36mm, height=38mm, xmin=-0.3, xmax=0.3, ymin=-1.3, ymax=1.3, xtick={-0.2,0,0.2}, xticklabels={$-0.2$,0,0.2}, ytick={-1,0,1}, yticklabels={$-eN_{\mathrm a}$,0,$eN_{\mathrm d}$}, xlabel={$x$}, ylabel={$\rho_{\mathrm q}$}]
  \draw[bkBlue, line width=1.4pt] (axis cs:-0.3,0) -- (axis cs:-0.215,0) -- (axis cs:-0.215,-1) -- (axis cs:0,-1) -- (axis cs:0,1) -- (axis cs:0.215,1) -- (axis cs:0.215,0) -- (axis cs:0.3,0);
\end{axis}
\begin{axis}[bkplot, axis lines=left, name=ef, at={(rq.east)}, anchor=west, xshift=22mm, width=36mm, height=38mm, xmin=-0.3, xmax=0.3, ymin=-3.8, ymax=0.4, xtick={-0.2,0,0.2}, xticklabels={$-0.2$,0,0.2}, ytick={-3.33,0}, yticklabels={$-3.3$,0}, xlabel={$x$}, ylabel={$\mathcal E$（$10^4$ V/cm）}]
  \draw[bkOrange, line width=1.4pt] (axis cs:-0.3,0) -- (axis cs:-0.215,0) -- (axis cs:0,-3.33) -- (axis cs:0.215,0) -- (axis cs:0.3,0);
\end{axis}
\begin{axis}[bkplot, axis lines=left, at={(ef.east)}, anchor=west, xshift=42mm, width=36mm, height=38mm, xmin=-0.3, xmax=0.3, ymin=-0.05, ymax=0.85, xtick={-0.2,0,0.2}, xticklabels={$-0.2$,0,0.2}, ytick={0,0.716}, yticklabels={0,0.72}, xlabel={$x$}, ylabel={$\phi$（V）}]
  \addplot[bkcurve3, line width=1.4pt, domain=-0.215:0, samples=30] {0.358*((x+0.215)/0.215)^2};
  \addplot[bkcurve3, line width=1.4pt, domain=0:0.215, samples=30] {0.716-0.358*((0.215-x)/0.215)^2};
  \draw[bkGreen, line width=1.4pt] (axis cs:-0.3,0) -- (axis cs:-0.215,0) (axis cs:0.215,0.716) -- (axis cs:0.3,0.716);
\end{axis}
\end{tikzpicture}
\bkfigcaption{例題7 の pn 接合（$N_{\mathrm a}=N_{\mathrm d}=\qty{1e16}{cm^{-3}}$。横軸は位置 $x$〔\si{\micro m}〕）。左：空乏層の電荷密度。中：電場（負は $x$ と反対の向き、つまり n形から p形へ）。右：電位。電荷を1回積分すると電場、もう1回積分すると電位になる。}\label{ph:原-9:fig-depl}
\end{bkfigure}

\phboxhead{ショックレーの式を導く}
順方向の電圧 $V$ で坂が $eV$ 下がると、ボルツマン因子から、n形の空乏層のはしにいる正孔（少数キャリア）は、つり合いのときの $p_{\mathrm n0}=n_{\mathrm i}^{\,2}/N_{\mathrm d}$ の $e^{eV/(k_{\mathrm B}T)}$ 倍になります。増えた分 $\Delta p=p_{\mathrm n0}(e^{eV/(k_{\mathrm B}T)}-1)$ は、n形の奥へ拡散しながら、寿命 $\tau_{\mathrm p}$ で消えていきます。「拡散で入ってくる分＝消える分」の式（数学 8-13 の熱方程式に、消える項を足した形）
\[ D_{\mathrm p}\frac{\dd^2\Delta p}{\dd x^2}=\frac{\Delta p}{\tau_{\mathrm p}} \]
の解は、$\Delta p(x)=\Delta p(0)\,e^{-x/L_{\mathrm p}}$（$L_{\mathrm p}=\sqrt{D_{\mathrm p}\tau_{\mathrm p}}$。数学 6-12 の特性方程式）です。はしでの拡散の電流の密度は $eD_{\mathrm p}\,\Delta p(0)/L_{\mathrm p}$。p形へ入る電子も同じように数えて足し、面積 $S$ をかけると、
\pheq{I=eS\,n_{\mathrm i}^{\,2}\Bigl(\frac{D_{\mathrm p}}{L_{\mathrm p}N_{\mathrm d}}+\frac{D_{\mathrm n}}{L_{\mathrm n}N_{\mathrm a}}\Bigr)\left(e^{eV/(k_{\mathrm B}T)}-1\right)}{A}
\label{ph:原-9:shockley}となります。かっこの前が⑦の飽和電流 $I_{\mathrm s}$ です（ショックレー、1949年）。

\begin{phflow}[拡散からダイオードの式へ]
\phflowrow{$D_{\mathrm p}\dfrac{\dd^2\Delta p}{\dd x^2}=\dfrac{\Delta p}{\tau_{\mathrm p}}$}{はしの数を決め、\\はしの流れを数える}{$I=I_{\mathrm s}\bigl(e^{eV/(k_{\mathrm B}T)}-1\bigr)$}
\end{phflow}

\noindent たとえば、p形を n形よりずっと濃くした接合（p$^+$n）で、$S=\qty{1}{mm^2}$、$N_{\mathrm d}=\qty{1e16}{cm^{-3}}$、$\tau_{\mathrm p}=\qty{1}{\micro s}$ とします。
\[ D_{\mathrm p}=480\times0.0259\approx\qty{12.4}{cm^2/s},\qquad L_{\mathrm p}=\sqrt{12.4\times10^{-6}}\approx\qty{3.5e-3}{cm} \]
p形が濃いので、2つめの項は小さく、無視できます。
\[ I_{\mathrm s}=1.60\times10^{-19}\times0.01\times(10^{10})^2\times\frac{12.4}{3.5\times10^{-3}\times10^{16}}\approx\qty{5.6e-14}{A} \]
です。⑦で使った $\qty{1e-14}{A}$ と、けたが合います。$I_{\mathrm s}$ は $n_{\mathrm i}^{\,2}\propto e^{-E_{\mathrm g}/(k_{\mathrm B}T)}$ に比例するので、温度が上がると急に大きくなります。そのため、同じ電流を流すのに要る電圧は、温度が 1 K 上がるごとに約 2 mV 下がります（体温計や IC の温度センサーは、これを使います）。

\begin{senmonsho}[半導体デバイスの専門書の書き方]
\begin{bkhonbun}
pn 接合ダイオードの電流密度は，理想的な場合
\[ J=J_0\left[\exp\left(\frac{qV}{nkT}\right)-1\right],\quad J_0=\frac{qD_pp_{n0}}{L_p}+\frac{qD_nn_{p0}}{L_n} \]
で与えられる．ここで $n$ は理想係数であり，拡散電流が支配的なとき $n=1$，空乏層内の再結合電流が支配的なとき $n=2$ に近づく．
\end{bkhonbun}
\bkyaku{専門書は、電気素量を $q$、ボルツマン定数を $k$ と書くことが多い（この本の $e$、$k_{\mathrm B}$）。$J$ は電流密度（この本の $j$）、$J_0$ は飽和電流の密度。$p_{n0}$ は「n形の中の、つり合いのときの正孔の数」（$=n_{\mathrm i}^{\,2}/N_{\mathrm d}$）。$\exp(x)$ は $e^x$ のこと。理想係数 $n$ は、⑦のモデル化の札の「10倍の電圧が 60 mV より大きくなる」ことを表す数で、キャリアの数 $n$ とは別のもの。}
\end{senmonsho}

\Needspace{8\baselineskip}\bksubsection{13}{ペルチェ素子とゼーベック効果}
\label{ph:原-9:sec13}
\phboxhead{温度の差から電圧：ゼーベック効果}
n形の半導体の棒の片はしを温めると、熱い側の電子のほうが、エネルギーが大きく、数も多くなるので、電子は熱い側から冷たい側へ拡散します（⑤）。冷たい側に電子がたまって－になり、熱い側は＋になって、電圧が生まれます。これを\textbf{ゼーベック効果}といいます。\bkidx{ぜーべっくこうか}{ゼーベック効果}p形なら、正孔が冷たい側へ移るので、向きが逆になります。

物理シリーズの『電磁気学』（電-4 ⑩）の熱電対は、2種類の金属で同じことをしていました。K 熱電対で 1 K あたり約 \qty{41}{\micro V} です。金属では、電気を運ぶ電子のほとんどが、フェルミ準位（$E_{\mathrm F}$）のすぐ近くにいて（原-8 ⑤）、温度でエネルギーがあまり変わらないので、電圧は小さいのです。半導体の電子は、$E_{\mathrm F}$ より $E_{\mathrm c}-E_{\mathrm F}$（0.2 eV ほど）も上にいて、1個が大きなエネルギーを運ぶので、1 K あたり数百 \si{\micro V} の電圧が出ます。

\phboxhead{電流で熱を運ぶ：ペルチェ効果}
反対に、n形の棒に電流を流すと、電子は、自分の運ぶエネルギーといっしょに、－極の側から＋極の側へ移ります。電子が入ってくる側（－極の側）の電極では、電子が金属から半導体の伝導帯へ「坂をのぼる」ので、周りから熱をうばって冷え、出ていく側では熱を出します。これが\textbf{ペルチェ効果}です。\bkidx{ぺるちぇこうか}{ペルチェ効果}n形と p形の小さな柱を、電気の回路としては交互に直列に、熱の通り道としては並列につなぐと、片面が冷え、反対の面が熱くなる板ができます（図\ref{ph:原-9:fig-peltier}）。これを\textbf{ペルチェ素子}といいます。

\begin{bkfigure}
\begin{tikzpicture}[x=1mm, y=1mm]
  \fill[bkBlue!12] (0,24) rectangle (88,29);
  \draw[bkInk, line width=0.6pt] (0,24) rectangle (88,29);
  \node[txj lab, anchor=south] at (44,29.5) {冷える面（熱を吸う）};
  \fill[bkRed!12] (0,0) rectangle (88,5);
  \draw[bkInk, line width=0.6pt] (0,0) rectangle (88,5);
  \node[txj lab, anchor=north west] at (0,-0.5) {熱くなる面（熱を出す）};
  \fill[bkGray!60] (8,22) rectangle (38,24) (48,22) rectangle (78,24);
  \fill[bkGray!60] (2,5) rectangle (18,7) (28,5) rectangle (58,7) (68,5) rectangle (84,7);
  \foreach \xx/\t/\cc in {8/p/bkRed!20, 28/n/bkBlue!25, 48/p/bkRed!20, 68/n/bkBlue!25} {
    \filldraw[fill=\cc, draw=bkInk, line width=0.6pt] (\xx,7) rectangle (\xx+10,22);
    \node[font=\sffamily\bfseries\footnotesize] at (\xx+5,17.5) {\t};}
  \foreach \xx in {8,48} \draw[bkGreen!50!black, line width=0.9pt, -{Stealth[length=2mm]}] (\xx+5,14.5) -- (\xx+5,9);
  \foreach \xx in {28,68} \draw[bkGreen!50!black, line width=0.9pt, -{Stealth[length=2mm]}] (\xx+5,9) -- (\xx+5,14.5);
  \draw[bkInk, line width=0.7pt] (2,6) -- (-6,6) -- (-6,-9) -- (94,-9) -- (94,6) -- (84,6);
  \txjE{(44,-9)}{180}
  \node[txj lab, anchor=north] at (44,-13.5) {直流の電源};
  \draw[bkRed!75!black, line width=1.6pt, -{Stealth[length=2.6mm]}] (100,27) -- (100,2);
  \node[txj lab, text=bkRed!75!black, anchor=west, align=left] at (102,14.5) {熱が\\運ばれる\\向き};
\end{tikzpicture}
\bkfigcaption{ペルチェ素子のつくり（横から見た図）。p形と n形の柱を、上と下の金属の板（灰色）で交互に直列につなぐ。柱の中の矢印は電流の向き。p形では正孔が電流と同じ向きに、n形では電子が電流と反対向きに動くので、どの柱でもキャリアは上から下へ動き、熱を上の面から下の面へ運ぶ。電流を逆にすると、冷える面と熱くなる面が入れかわる。}\label{ph:原-9:fig-peltier}
\end{bkfigure}

\noindent ペルチェ素子は、気体も、動く部品も使わないヒートポンプです。小さな冷蔵庫、カメラのセンサーやレーザーの温度を一定に保つ装置、コンピュータの冷却などに使われます。電流の向きを逆にすると、冷やす面と温める面が入れかわります。

\phboxhead{ペルチェ素子も、第二法則にしたがう}
物理シリーズの『熱学』（熱-5）では、ペルチェ素子のような気体を使わないヒートポンプも、第二法則の上限にしたがうと書きました。熱を低温 $T_{\mathrm L}$ から高温 $T_{\mathrm H}$ へくみ上げる機械の成績係数は、どんなしくみでも、カルノーの値（熱-5 ⑤）
\pheq{\mathrm{COP}_{\text{冷}}\leqq\frac{T_{\mathrm L}}{T_{\mathrm H}-T_{\mathrm L}}}{1 ＝ K ÷ K}
\label{ph:原-9:cop}をこえられません。実際のペルチェ素子では、柱を流れる電流のジュール熱と、熱い面から冷たい面へもどる熱伝導が大きな損失になり、成績係数は、ふつう1より小さい値です。それでも、小さく、静かで、温度を細かく調整できるので、使い道が広いのです。

\phboxhead{ゼーベック効果で発電する}
ペルチェ素子の両面に温度の差をつけると、ゼーベック効果で電圧が出て、発電機になります。たとえば、1対（p と n の柱1組）で 1 K あたり \qty{0.4}{mV} の素子が127対あれば、30 K の差で $0.4\times10^{-3}\times127\times30\approx\qty{1.5}{V}$ です（K 熱電対の約10倍の電圧が、1対で出ます）。工場や車の捨てる熱から電気をとり出す研究が進められています。

%% ======================== まだ・この先 ========================
\begin{phmada}[青色 LED は、なぜ大きな電流で効率が下がる？]
窒化ガリウムの青色 LED は、小さな電流では、流しこんだ電子の大部分を光子に変えます。ところが、電流を大きくして明るくしようとすると、電子1個あたりの光の量（効率）が下がってしまいます。これを\textbf{ドループ}（垂れ下がり）といいます。照明を明るくするには、LED の数をふやすか、大きな電流で効率が下がるのをがまんするしかなく、大きな問題です。

原因として、電子と正孔が出会ったとき、光を出さずに、3つめのキャリアにエネルギーを渡してしまう「オージェ再結合」、電子が p形の側へもれ出すこと、結晶の欠陥などがあげられてきました。2013年には、LED から飛び出す電子のエネルギーをはかって、オージェ再結合が主な原因だとする実験が発表されました（アイブランドら）。一方で、もれや欠陥の寄与も大きいという研究もあり、材料や構造によって、どれがどれだけ効いているのかは、議論が続いています。
\end{phmada}

\begin{phsaki}{半導体の先へ}
この章の pn 接合は、半導体工学・電子デバイス工学の出発点です。専門の本では、禁制帯の違う半導体を原子の層ごとに重ねた接合（ヘテロ接合）で、電子と正孔を薄い層に閉じこめて光らせる半導体レーザーや、何段も重ねた高い効率の太陽電池を学びます（ヘテロ接合の研究は 2000年のノーベル物理学賞）。太陽光発電工学では、パネルを何枚もつないで、天気で変わる光から、いちばん大きな電力をとり出し続ける制御（最大電力点の追従）を学び、インバーター（電-16）で交流にして送ります。トランジスタ（原-10）も、この章の pn 接合を2つ、または電場で作る接合を使っています。
\end{phsaki}


\begin{taisetsu}
\def\tsl#1{\makebox[7.2\zw][l]{\sffamily\bfseries\small\color{\bkTT}#1}}
\tsl{キャリアの数}$n_{\mathrm i}=\sqrt{N_{\mathrm c}N_{\mathrm v}}\,e^{-E_{\mathrm g}/(2k_{\mathrm B}T)}$、$np=n_{\mathrm i}^{\,2}$、n形 $n\approx N_{\mathrm d}$、p形 $p\approx N_{\mathrm a}$\\[3pt]
\tsl{流れ}$\sigma=e(n\mu_{\mathrm n}+p\mu_{\mathrm p})$、$D=\mu k_{\mathrm B}T/e$、$L=\sqrt{D\tau}$\\[3pt]
\tsl{pn 接合}$V_{\mathrm{bi}}=\dfrac{k_{\mathrm B}T}{e}\ln\dfrac{N_{\mathrm a}N_{\mathrm d}}{n_{\mathrm i}^{\,2}}$、$W=\sqrt{\dfrac{2\varepsilon V_{\mathrm{bi}}}{e}\Bigl(\dfrac1{N_{\mathrm a}}+\dfrac1{N_{\mathrm d}}\Bigr)}$\\[3pt]
\tsl{ダイオード}$I=I_{\mathrm s}(e^{eV/(k_{\mathrm B}T)}-1)$、約 60 mV で電流10倍（300 K）\\[3pt]
\tsl{LED・太陽電池}$\lambda\approx hc/E_{\mathrm g}$、$I=I_{\mathrm L}-I_{\mathrm s}(e^{eV/(k_{\mathrm B}T)}-1)$、$V_{\mathrm{oc}}=\dfrac{k_{\mathrm B}T}{e}\ln\Bigl(\dfrac{I_{\mathrm L}}{I_{\mathrm s}}+1\Bigr)$\\[3pt]
\tsl{光をはかる}感度 $=\eta e\lambda/(hc)$\\[3pt]
\tsl{ペルチェ素子}$\mathrm{COP}_{\text{冷}}\leqq T_{\mathrm L}/(T_{\mathrm H}-T_{\mathrm L})$
\end{taisetsu}

\bkkigocards{\bkkigocard{$N_{\mathrm d}$}{エヌ・ディー}{ドナーの数密度〔cm$^{-3}$〕。アクセプターは $N_{\mathrm a}$。アボガドロ定数 $N_{\mathrm A}$ とは別の量}{donor density}{原-9}}
{\bkkigocard{$\mathcal E$}{イー（筆記体）}{原-9 では半導体の中の電場〔V/cm〕。第4巻の起電力 $\mathcal E$ とは別の量}{electric field}{原-9}}
\bkkigocards{\bkkigocard{$\mu_{\mathrm n}$}{ミュー・エヌ}{電子の移動度〔cm$^2$/(V·s)〕。正孔は $\mu_{\mathrm p}$。第4巻の透磁率 $\mu$ とは別の量}{mobility}{原-9}}{}

\begin{kotoba}
\small
\bkword{正孔}{結合の手から電子が抜けた穴。＋の電荷 $e$ をもつ粒子のように動く。}
\bkword{n形・p形}{ドナー（リンなど）を混ぜて電子を多くした半導体と、アクセプター（ホウ素など）を混ぜて正孔を多くした半導体。}
\bkword{真性キャリア密度}{まじりけのない半導体の電子（＝正孔）の数密度 $n_{\mathrm i}$。}
\bkword{移動度・拡散係数}{電場1 V/cm あたりの平均の速さ $\mu$ と、数の坂で広がる速さの係数 $D$。$D=\mu k_{\mathrm B}T/e$。}
\bkword{空乏層・内蔵電位}{pn 接合の境目の、キャリアがいない層と、その両側の電位の差 $V_{\mathrm{bi}}$。}
\bkword{ダイオード}{pn 接合で、電流を順方向にだけ通す部品。}
\bkword{開放電圧・曲線因子}{太陽電池の、何もつながないときの電圧と、$P_{\max}/(V_{\mathrm{oc}}I_{\mathrm L})$。}
\bkword{ゼーベック効果・ペルチェ効果}{温度の差で電圧が生まれることと、電流で熱が運ばれること。}
\end{kotoba}

%% ======================== 例題 ========================
\Needspace{10\baselineskip}\bksection{reidai}{例題でたしかめよう}
\label{ph:原-9:ex1}
\begin{phreidai}[LED の色（はてな？）]
赤い LED の材料の禁制帯幅を約 \qty{1.9}{eV}、青い LED の材料を約 \qty{2.7}{eV} とします。それぞれの LED が出す光の波長を求めましょう。
\phbamen
LED の接合で、電子が伝導帯から価電子帯へ落ち、光子を1個出す場面です。
\phmushi
電子が伝導帯の底より少し上から落ちる分（熱の $k_{\mathrm B}T$ くらい）は無視し、光子のエネルギーは $E_{\mathrm g}$ と等しいとします。
\phzu
$E_{\mathrm g}=\qty{1.9}{eV}$（赤）・\qty{2.7}{eV}（青）、$hc=\qty{1240}{eV.nm}$、波長 $\lambda$。
\phriyu
光子のエネルギーと波長の関係 $E=hc/\lambda$（原-2）を、$E=E_{\mathrm g}$ として使います（⑨）。
\phtateru
\[ \lambda=\frac{hc}{E_{\mathrm g}}=\frac{1240\ \text{eV·nm}}{E_{\mathrm g}} \]
\phtoku
\[ \text{赤}\ \ \lambda=\frac{1240}{1.9}\approx\qty{650}{nm},\qquad \text{青}\ \ \lambda=\frac{1240}{2.7}\approx\qty{460}{nm} \]
\phtashikame
単位：eV·nm÷eV＝nm。\\
大きさ：目に見える光は約 380〜780 nm。赤は長いほう、青は短いほうで、色の並びに合います。\\
逆向き：650 nm の光子のエネルギーは $1240/650\approx\qty{1.9}{eV}$ で、もとにもどります。
\kotae 赤 約\qty{650}{nm}、青 約\qty{460}{nm}
\end{phreidai}
\begin{jibun}
禁制帯幅が約 \qty{2.3}{eV} の緑の LED と、\qty{1.42}{eV} のガリウムヒ素の LED の波長を求めましょう。ガリウムヒ素の光は目に見えますか。
\bkjans{緑 $1240\div2.3\approx\qty{540}{nm}$、ガリウムヒ素 $1240\div1.42\approx\qty{870}{nm}$（赤外線で、目に見えない）}
\end{jibun}

\label{ph:原-9:ex2}
\begin{phreidai}[ケイ素の電子の数と温度（サーミスター）]
ケイ素の $N_{\mathrm c}=\qty{2.8e19}{cm^{-3}}$、$N_{\mathrm v}=\qty{1.04e19}{cm^{-3}}$、$E_{\mathrm g}=\qty{1.12}{eV}$ として、300 K の真性キャリア密度 $n_{\mathrm i}$ を求めましょう。また、350 K では何倍になりますか（$N_{\mathrm c}N_{\mathrm v}$ は $T^3$ に比例し、$E_{\mathrm g}$ は変わらないとします）。
\phbamen
まじりけのないケイ素の中で、熱で結合の手から外れる電子の数を数えます。電子体温計のサーミスターが温度を感じるしくみです。
\phmushi
$E_{\mathrm g}$ が温度で少し変わることは無視します。
\phzu
$T=\qty{300}{K}$・\qty{350}{K}。温度を変えるので、$k_{\mathrm B}T$ は $\qty{8.617e-5}{eV/K}\times T$ で計算します（300 K で 0.02585 eV、350 K で 0.03016 eV）。
\phriyu
③の $n_{\mathrm i}=\sqrt{N_{\mathrm c}N_{\mathrm v}}\,e^{-E_{\mathrm g}/(2k_{\mathrm B}T)}$ を使います。
\phtateru
\[ n_{\mathrm i}(300)=\sqrt{2.8\times10^{19}\times1.04\times10^{19}}\times e^{-1.12/(2\times0.02585)} \]
\[ \frac{n_{\mathrm i}(350)}{n_{\mathrm i}(300)}=\Bigl(\frac{350}{300}\Bigr)^{3/2}\times\frac{e^{-18.57}}{e^{-21.66}}=\Bigl(\frac{350}{300}\Bigr)^{3/2}e^{\,21.66-18.57} \]
\phtoku
\[ \sqrt{2.91\times10^{38}}=1.71\times10^{19},\qquad e^{-21.66}=3.9\times10^{-10} \]
なので、$n_{\mathrm i}(300)\approx\qty{6.7e9}{cm^{-3}}$。比は $1.26\times e^{3.10}=1.26\times22.1\approx28$ 倍。
\phtashikame
単位：cm$^{-3}$×（数）＝cm$^{-3}$。\\
大きさ：教科書の $\qty{1.0e10}{cm^{-3}}$ と同じけたです（③で見たとおり、肩の数のわずかな違いで値が変わる）。\\
極端な場合：温度を0に近づけると、肩が $-\infty$ になり $n_{\mathrm i}\to0$。冷やすと、ケイ素は絶縁体に近づきます。
\kotae $n_{\mathrm i}\approx\qty{6.7e9}{cm^{-3}}$、350 K で約28倍
\end{phreidai}
\begin{jibun}
例題2 と同じケイ素で、300 K から 330 K に上げると、$n_{\mathrm i}$ は何倍になりますか（330 K の $k_{\mathrm B}T=\qty{0.02844}{eV}$）。
\bkjans{$(330/300)^{3/2}\times e^{21.66-19.69}=1.154\times7.2\approx8.3$ 倍（約8倍）}
\end{jibun}

\label{ph:原-9:ex3}
\begin{phreidai}[リンを混ぜたケイ素の電子と正孔]
ケイ素に、リンを $N_{\mathrm d}=\qty{1e16}{cm^{-3}}$ だけ混ぜました。300 K の電子と正孔の数密度を求めましょう。$n_{\mathrm i}=\qty{1.0e10}{cm^{-3}}$ とします。
\phbamen
n形の半導体の中のキャリアを数える場面です。
\phmushi
リンはすべて電子を出しているとします。熱で結合の手から外れる電子は、リンからの電子に比べて無視します。
\phzu
$N_{\mathrm d}=10^{16}$、$n_{\mathrm i}=10^{10}$（cm$^{-3}$）。
\phriyu
多いほうの電子は $n\approx N_{\mathrm d}$、少ないほうの正孔は $np=n_{\mathrm i}^{\,2}$（④）から求めます。
\phtateru
\[ n=N_{\mathrm d},\qquad p=\frac{n_{\mathrm i}^{\,2}}{n} \]
\phtoku
\[ n=\qty{1e16}{cm^{-3}},\qquad p=\frac{(10^{10})^2}{10^{16}}=\qty{1e4}{cm^{-3}} \]
\phtashikame
単位：cm$^{-6}$÷cm$^{-3}$＝cm$^{-3}$。\\
大きさ：ケイ素の原子 $\num{5.0e22}$ 個に対してリンは $10^{16}$ 個、つまり500万個に1個。それだけで電子は $n_{\mathrm i}$ の100万倍になります。\\
確かめ：$np=10^{20}=n_{\mathrm i}^{\,2}$。仮定どおり、$n_{\mathrm i}$ は $N_{\mathrm d}$ よりずっと小さい。
\kotae 電子 $\qty{1e16}{cm^{-3}}$、正孔 $\qty{1e4}{cm^{-3}}$
\end{phreidai}
\begin{jibun}
ケイ素にホウ素を $N_{\mathrm a}=\qty{2e15}{cm^{-3}}$ 混ぜました。正孔と電子の数密度を求めましょう（$n_{\mathrm i}=\qty{1.0e10}{cm^{-3}}$）。
\bkjans{正孔 $p\approx\qty{2e15}{cm^{-3}}$、電子 $n=10^{20}\div(2\times10^{15})=\qty{5e4}{cm^{-3}}$}
\end{jibun}

\label{ph:原-9:ex4}
\begin{phreidai}[混ぜると抵抗率が何けたも変わる]
例題3 の n形のケイ素と、まじりけのないケイ素（$n=p=\qty{1.0e10}{cm^{-3}}$）の抵抗率を求めましょう。移動度は電子 1350、正孔 480（cm$^2$/(V·s)）とします。
\phbamen
同じケイ素の板で、まじりけのあるなしで、電気の通しやすさを比べます。
\phmushi
移動度は、まじりけの量によらないとします（実際は、混ぜる量が多いと少し下がる）。
\phzu
$e=\qty{1.60e-19}{C}$、例題3 から $n=10^{16}$、$p=10^4$。
\phriyu
導電率は、電子と正孔の両方の流れを足した $\sigma=e(n\mu_{\mathrm n}+p\mu_{\mathrm p})$（⑤）、抵抗率は $\rho=1/\sigma$ です。
\phtateru
\[ \rho_{\mathrm n形}=\frac{1}{e\,n\mu_{\mathrm n}},\qquad \rho_{\text{まじりけなし}}=\frac{1}{e\,n_{\mathrm i}(\mu_{\mathrm n}+\mu_{\mathrm p})} \]
\phtoku
n形：$\sigma=1.60\times10^{-19}\times10^{16}\times1350=\qty{2.16}{S/cm}$ なので、$\rho\approx\qty{0.46}{\ohm.cm}$（正孔の分は $10^{12}$ 分の1で無視）。\\
まじりけなし：$\sigma=1.60\times10^{-19}\times10^{10}\times1830=\qty{2.93e-6}{S/cm}$ なので、$\rho\approx\qty{3.4e5}{\ohm.cm}$。
\phtashikame
単位：C×cm$^{-3}$×cm$^2$/(V·s)＝A/(V·cm)＝S/cm。\\
大きさ：抵抗率の比は約 $7\times10^5$。混ぜる量で、抵抗率を何けたも自由に変えられます。\\
比べる：銅の抵抗率は約 $\qty{1.7e-6}{\ohm.cm}$（電-4）。n形のケイ素は、それより約27万倍大きい。
\kotae n形 約\qty{0.46}{\ohm.cm}、まじりけなし 約 $\qty{3.4e5}{\ohm.cm}$
\end{phreidai}
\begin{jibun}
ホウ素を $N_{\mathrm a}=\qty{2e15}{cm^{-3}}$ 混ぜた p形のケイ素の抵抗率を求めましょう（正孔の移動度 480 cm$^2$/(V·s)）。
\bkjans{$\sigma=1.60\times10^{-19}$ $\times2\times10^{15}\times480$ $\approx\qty{0.154}{S/cm}$、抵抗率 約\qty{6.5}{\ohm.cm}}
\end{jibun}

\label{ph:原-9:ex5}
\begin{phreidai}[n形のフェルミ準位の位置]
例題3 の n形のケイ素（$N_{\mathrm d}=\qty{1e16}{cm^{-3}}$）のフェルミ準位は、伝導帯の下のはしから何 eV 下にありますか（$N_{\mathrm c}=\qty{2.8e19}{cm^{-3}}$、300 K）。
\phbamen
帯の図（図\ref{ph:原-9:fig-pn}）で、n形の $E_{\mathrm F}$ をかき入れるときの位置を求めます。
\phmushi
フェルミ分布の $+1$ を無視するボルツマンの裾の近似（⑪）を使います。
\phzu
$n=N_{\mathrm d}=10^{16}$、$k_{\mathrm B}T=\qty{0.0259}{eV}$。
\phriyu
⑪の $n=N_{\mathrm c}e^{-(E_{\mathrm c}-E_{\mathrm F})/(k_{\mathrm B}T)}$ を $E_{\mathrm c}-E_{\mathrm F}$ について解きます。
\phtateru
\[ E_{\mathrm c}-E_{\mathrm F}=k_{\mathrm B}T\ln\frac{N_{\mathrm c}}{n} \]
\phtoku
\[ E_{\mathrm c}-E_{\mathrm F}=0.0259\times\ln\frac{2.8\times10^{19}}{10^{16}}=0.0259\times7.94\approx\qty{0.21}{eV} \]
\phtashikame
単位：eV×（数）＝eV。\\
仮定の確かめ：$0.21$ eV は $k_{\mathrm B}T$ の約8倍なので、$+1$ を無視してよい（$e^{8}\approx3000\gg1$）。\\
極端な場合：$N_{\mathrm d}$ を $N_{\mathrm c}$ に近づけると、$E_{\mathrm F}$ は伝導帯のはしに近づきます。
\kotae 伝導帯の下のはしより約\qty{0.21}{eV}下
\end{phreidai}
\begin{jibun}
ホウ素を $N_{\mathrm a}=\qty{1e16}{cm^{-3}}$ 混ぜた p形のケイ素のフェルミ準位は、価電子帯の上のはしから何 eV 上にありますか（$N_{\mathrm v}=\qty{1.04e19}{cm^{-3}}$）。
\bkjans{$E_{\mathrm F}-E_{\mathrm v}=0.0259\times\ln(1.04\times10^{19}\div10^{16})=0.0259\times6.95\approx\qty{0.18}{eV}$}
\end{jibun}

\label{ph:原-9:ex6}
\begin{phreidai}[電子の拡散係数と拡散距離]
p形のケイ素の中に入った電子（移動度 \qty{1350}{cm^2/(V.s)}）の拡散係数を求めましょう。寿命が \qty{1}{\micro s} なら、電子は消えるまでに平均どれくらい進みますか（300 K）。
\phbamen
太陽電池の p形の層で、光で生まれた電子が接合まで届くかを調べる場面です。
\phmushi
電場のない所での拡散だけを考えます。
\phzu
$\mu_{\mathrm n}=1350$、$k_{\mathrm B}T/e=\qty{0.0259}{V}$、$\tau_{\mathrm n}=\qty{1e-6}{s}$。
\phriyu
アインシュタインの関係 $D=\mu k_{\mathrm B}T/e$ と、拡散距離 $L=\sqrt{D\tau}$（⑤）を使います。
\phtateru
\[ D_{\mathrm n}=\mu_{\mathrm n}\frac{k_{\mathrm B}T}{e},\qquad L_{\mathrm n}=\sqrt{D_{\mathrm n}\tau_{\mathrm n}} \]
\phtoku
\[ D_{\mathrm n}=1350\times0.0259\approx\qty{35}{cm^2/s} \]
\[ L_{\mathrm n}=\sqrt{35\times10^{-6}}\approx\qty{5.9e-3}{cm}=\qty{59}{\micro m} \]
\phtashikame
単位：cm$^2$/(V·s)×V＝cm$^2$/s。$\sqrt{\text{cm}^2\text{/s}\times\text{s}}$＝cm。\\
大きさ：髪の毛の太さくらい。空乏層（\qty{1}{\micro m} ほど）よりずっと長いので、空乏層の外で生まれた電子も接合に届きます。\\
2倍の確かめ：寿命を4倍にすると、$L$ は2倍。
\kotae $D_{\mathrm n}\approx\qty{35}{cm^2/s}$、$L_{\mathrm n}\approx\qty{59}{\micro m}$
\end{phreidai}
\begin{jibun}
n形のケイ素の中に入った正孔（移動度 \qty{480}{cm^2/(V.s)}、寿命 \qty{1}{\micro s}）の拡散係数と拡散距離を求めましょう。
\bkjans{$D_{\mathrm p}=480\times0.0259\approx\qty{12}{cm^2/s}$、$L_{\mathrm p}=\sqrt{12.4\times10^{-6}}\approx\qty{3.5e-3}{cm}=\qty{35}{\micro m}$}
\end{jibun}

\label{ph:原-9:ex7}
\begin{phreidai}[pn 接合の内蔵電位・空乏層・電場]
どちらも $\qty{1e16}{cm^{-3}}$ の p形と n形でできたケイ素の pn 接合があります。300 K の内蔵電位、空乏層の厚さ、いちばん強い電場を求めましょう（$n_{\mathrm i}=\qty{1.0e10}{cm^{-3}}$、比誘電率 11.7）。
\phbamen
電圧をかけていない pn 接合の、境目のようすを数で調べます。
\phmushi
境目で p形から n形へ急に変わる（段の形）とし、空乏層の中にキャリアはいない、外に電場はないとします（空乏近似）。
\phzu
$N_{\mathrm a}=N_{\mathrm d}=10^{16}$、$\varepsilon=11.7\times\qty{8.85e-14}{F/cm}$、$e=\qty{1.60e-19}{C}$。
\phriyu
$V_{\mathrm{bi}}$ は⑥（導き方は⑪）、$W$ と $\mathcal E_{\max}$ は⑫のポアソン方程式から出した式を使います。
\phtateru
\[ V_{\mathrm{bi}}=\frac{k_{\mathrm B}T}{e}\ln\frac{N_{\mathrm a}N_{\mathrm d}}{n_{\mathrm i}^{\,2}},\qquad W=\sqrt{\frac{2\varepsilon V_{\mathrm{bi}}}{e}\Bigl(\frac1{N_{\mathrm a}}+\frac1{N_{\mathrm d}}\Bigr)} \]
\[ |\mathcal E_{\max}|=\frac{2V_{\mathrm{bi}}}{W} \]
\phtoku
$V_{\mathrm{bi}}=0.0259\times\ln10^{12}=0.0259\times27.6\approx\qty{0.72}{V}$。\\
$\varepsilon=11.7\times8.85\times10^{-14}=\qty{1.035e-12}{F/cm}$、$1/N_{\mathrm a}+1/N_{\mathrm d}=\qty{2e-16}{cm^3}$ なので、
\[ \frac{2\varepsilon V_{\mathrm{bi}}}{e}=\frac{2\times1.035\times10^{-12}\times0.716}{1.60\times10^{-19}}=\qty{9.26e6}{cm^{-1}} \]
\[ W=\sqrt{9.26\times10^{6}\times2\times10^{-16}}=\sqrt{1.85\times10^{-9}}\approx\qty{4.3e-5}{cm} \]
つまり $W\approx\qty{0.43}{\micro m}$。\\
$|\mathcal E_{\max}|=2\times0.716\div(4.3\times10^{-5})\approx\qty{3.3e4}{V/cm}$。
\phtashikame
単位：$\sqrt{\text{F/cm}\times\text{V}\div\text{C}\times\text{cm}^3}$＝cm。V÷cm＝V/cm。\\
大きさ：空乏層は光の波長くらいの厚さ。電場は $\qty{3.3e6}{V/m}$ で、空気が放電する電場（約 $\qty{3e6}{V/m}$）と同じくらい。\\
対称：$N_{\mathrm a}=N_{\mathrm d}$ なので、空乏層は p形と n形に半分ずつ（\qty{0.22}{\micro m} ずつ）広がります。
\kotae $V_{\mathrm{bi}}\approx\qty{0.72}{V}$、$W\approx\qty{0.43}{\micro m}$、$|\mathcal E_{\max}|\approx\qty{3.3e4}{V/cm}$
\end{phreidai}
\begin{jibun}
p形を $\qty{1e17}{cm^{-3}}$、n形を $\qty{1e16}{cm^{-3}}$ にした接合の、内蔵電位・空乏層の厚さ・いちばん強い電場を求めましょう。空乏層は、どちらの側に多く広がりますか。
\bkjans{$V_{\mathrm{bi}}=0.0259\times\ln10^{13}\approx\qty{0.78}{V}$、$W\approx\qty{0.33}{\micro m}$、$|\mathcal E_{\max}|\approx\qty{4.7e4}{V/cm}$。$N_{\mathrm a}x_{\mathrm p}=N_{\mathrm d}x_{\mathrm n}$ から、うすい n形の側に約9割（\qty{0.30}{\micro m}）}
\end{jibun}

\label{ph:原-9:ex8}
\begin{phreidai}[ダイオードの電圧と「60 mV で10倍」]
飽和電流 $I_{\mathrm s}=\qty{1e-14}{A}$ のダイオードに、順方向に \qty{1}{mA} を流すには、何 V かければよいでしょう。電流を10倍の \qty{10}{mA} にするには、電圧をあと何 V 上げればよいでしょう（300 K）。
\phbamen
電子回路で、ダイオードにかかる電圧を見積もる場面です。数学 3-12 の「約 0.06 V ごとに電流10倍」を、式で確かめます。
\phmushi
直列抵抗は無視し、理想ダイオードの式を使います。
\phzu
$I=\qty{1e-3}{A}$、$I_{\mathrm s}=\qty{1e-14}{A}$、$k_{\mathrm B}T/e=\qty{0.0259}{V}$。
\phriyu
⑦のダイオードの式 $I=I_{\mathrm s}(e^{eV/(k_{\mathrm B}T)}-1)$ を $V$ について解きます。$I\gg I_{\mathrm s}$ なので $-1$ は無視できます。
\phtateru
\[ V=\frac{k_{\mathrm B}T}{e}\ln\frac{I}{I_{\mathrm s}},\qquad \Delta V=\frac{k_{\mathrm B}T}{e}\ln10 \]
\phtoku
\[ V=0.0259\times\ln10^{11}=0.0259\times25.3\approx\qty{0.66}{V} \]
\[ \Delta V=0.0259\times2.303\approx\qty{0.060}{V} \]
\phtashikame
単位：V×（数）＝V。\\
大きさ：ケイ素のダイオードの立ち上がり（0.6〜0.7 V）に合います。\\
3-12 との照合：$\Delta V=\qty{59.6}{mV}$ は「約 0.06 V」と同じ。$I_{\mathrm s}$ によらず、温度だけで決まります。
\kotae 約\qty{0.66}{V}。10倍にするには約\qty{0.06}{V}（\qty{60}{mV}）上げる
\end{phreidai}
\begin{jibun}
$I_{\mathrm s}=\qty{1e-13}{A}$ のダイオードに \qty{10}{mA} を流す電圧を求めましょう。電流を100倍にするには、電圧を何 V 上げればよいですか。
\bkjans{$0.0259\times\ln10^{11}\approx\qty{0.66}{V}$。100倍は $0.0259\times\ln100\approx\qty{0.12}{V}$}
\end{jibun}

\label{ph:原-9:ex9}
\begin{phreidai}[一方通行の度合い（整流）]
例題8 のダイオードに、順方向に \qty{0.5}{V} かけたときと、逆方向に \qty{0.5}{V} かけたときの、電流の大きさの比を求めましょう。
\phbamen
ダイオードが、どれくらい「一方通行」なのかを、数で比べます。
\phmushi
逆方向の降伏や、表面をもれる電流は無視します。
\phzu
順方向 $V=+0.5$、逆方向 $V=-0.5$（V）。
\phriyu
⑦の式に、両方の電圧を入れて比べます。
\phtateru
\[ \frac{|I(+0.5)|}{|I(-0.5)|}=\frac{e^{0.5/0.0259}-1}{1-e^{-0.5/0.0259}} \]
\phtoku
$0.5/0.0259=19.3$ なので、分子は $e^{19.3}-1\approx2.4\times10^8$、分母は $1-e^{-19.3}\approx1$。比は約 $2.4\times10^8$ 倍です。
\phtashikame
大きさ：順方向は $10^{-14}\times2.4\times10^8\approx\qty{2.4e-6}{A}$、逆方向は $\qty{1e-14}{A}$。\\
極端な場合：$V\to0$ では、どちらの向きも0に近づき、比は1に近づきます（小さな電圧では一方通行にならない）。
\kotae 約 $2.4\times10^8$ 倍
\end{phreidai}
\begin{jibun}
同じダイオードに、順方向と逆方向に \qty{0.6}{V} ずつかけたときの、電流の比を求めましょう。
\bkjans{$e^{0.6/0.0259}\approx1.2\times10^{10}$ 倍（順方向は約 \qty{0.12}{mA}）}
\end{jibun}

\label{ph:原-9:ex10}
\begin{phreidai}[太陽電池のいちばん大きな電力と効率]
面積 \qty{100}{cm^2} のケイ素の太陽電池のセルがあります。1 m$^2$ に \qty{1000}{W} の光を当てたときの光電流の密度を \qty{35}{mA/cm^2}、飽和電流の密度を $\qty{1e-12}{A/cm^2}$ とします。開放電圧、いちばん大きな電力とそのときの電圧、曲線因子、効率を求めましょう（300 K）。
\phbamen
屋根の太陽光パネルの1枚のセルで、取り出せる電力を求めます。
\phmushi
直列抵抗と、光で温まって温度が上がることは無視します。
\phzu
\[ I_{\mathrm L}=0.035\times100=\qty{3.5}{A},\qquad I_{\mathrm s}=10^{-12}\times100=\qty{1e-10}{A} \]
当たる光は $1000\times0.01=\qty{10}{W}$。
\phriyu
⑩の $I=I_{\mathrm L}-I_{\mathrm s}(e^{eV/(k_{\mathrm B}T)}-1)$ を使い、$P=VI$ のいちばん大きな所を、表計算で探します。
\phtateru
\[ V_{\mathrm{oc}}=\frac{k_{\mathrm B}T}{e}\ln\Bigl(\frac{I_{\mathrm L}}{I_{\mathrm s}}+1\Bigr),\qquad P(V)=V\Bigl[I_{\mathrm L}-I_{\mathrm s}\bigl(e^{V/0.0259}-1\bigr)\Bigr] \]
\phtoku
$V_{\mathrm{oc}}=0.0259\times\ln(3.5\times10^{10})=0.0259\times24.3\approx\qty{0.63}{V}$。電圧を 0.01 V ずつ変えて $P$ を計算すると、
\begin{bkcenter}
\begin{bktable}{colspec={Q[c,m,wd=16mm] Q[c,m,wd=14mm] Q[c,m,wd=14mm] Q[c,m,wd=14mm] Q[c,m,wd=14mm] Q[c,m,wd=14mm]}, rowsep=0pt, row{1}={bg=bkTheme!10}, cells={font=\small}}
$V$〔V〕 & 0.50 & 0.53 & 0.55 & 0.56 & 0.58 \\
$I$〔A〕 & 3.48 & 3.42 & 3.33 & 3.25 & 2.97 \\
$P$〔W〕 & 1.74 & 1.81 & 1.83 & 1.82 & 1.72 \\
\end{bktable}
\end{bkcenter}
で、$V\approx\qty{0.55}{V}$ のとき $P_{\max}\approx\qty{1.83}{W}$。曲線因子は $1.83\div(0.629\times3.5)\approx0.83$、効率は $1.83\div10\approx18\,\%$。
\phtashikame
単位：V×A＝W。\\
大きさ：セル1枚 0.6 V ほど、2 W 弱。パネルは60枚ほどを直列にして 30 V 以上にします。\\
確かめ：$P_{\max}$ の電圧は $V_{\mathrm{oc}}$ より \qty{0.08}{V} ほど低い（$k_{\mathrm B}T/e$ の約3倍）。$\dd P/\dd V=0$（数学 5-7）で解いても同じ答えです。
\kotae $V_{\mathrm{oc}}\approx\qty{0.63}{V}$、$P_{\max}\approx\qty{1.83}{W}$（約\qty{0.55}{V}）、曲線因子 約0.83、効率 約18\,\%
\end{phreidai}
\begin{jibun}
例題10 のセルを \qty{0.50}{V} で使うと、取り出せる電力は何 W で、$P_{\max}$ の何\,\% ですか。また、光電流の密度が \qty{30}{mA/cm^2} に下がると、開放電圧はいくらになりますか。
\bkjans{\qty{0.50}{V} で $I\approx\qty{3.48}{A}$、$P\approx\qty{1.74}{W}$（$P_{\max}$ の約95\,\%）。開放電圧 $0.0259\times\ln(3\times10^{10})\approx\qty{0.62}{V}$}
\end{jibun}

\label{ph:原-9:ex11}
\begin{phreidai}[光電池で LED を光らせる]
理科室の光電池1個は、日光で約 \qty{0.6}{V} の電圧を出すとします。光らせるのに約 \qty{2}{V} が要る LED を光電池で光らせるには、光電池を何個、どうつなげばよいでしょう。
\phbamen
小学校の光電池で、モーターではなく LED を光らせたい場面です。
\phmushi
光電池の電圧は、流す電流で少し下がることを無視します。
\phzu
1個 \qty{0.6}{V}、LED に要る電圧 \qty{2}{V}、個数 $N$。
\phriyu
光電池1個の電圧は、⑩の開放電圧で決まり、光を強くしてもほとんど上がりません。電池と同じく、直列につなぐと電圧が足し算になります（電-5 ②）。
\phtateru
\[ 0.6\times N\geqq2 \]
\phtoku
$N\geqq3.33\cdots$ なので、$N=4$。3個では $\qty{1.8}{V}$ で足りません。
\phtashikame
大きさ：4個で \qty{2.4}{V}。電流は光電池1個分（光の強さで決まる）です。\\
注意：光電池を並列にしても、電圧は上がらず、電流が増えるだけです（①の2倍の図）。
\kotae 4個を直列につなぐ
\end{phreidai}
\begin{jibun}
1個 \qty{0.5}{V} の光電池で、約 \qty{3}{V} が要る青い LED を光らせるには、何個を直列につなげばよいですか。
\bkjans{$0.5\times N\geqq3$ より6個}
\end{jibun}

\label{ph:原-9:ex12}
\begin{phreidai}[フォトダイオードの感度]
光ファイバー通信の受け手のフォトダイオードに、波長 \qty{850}{nm} の光が当たります。当たった光子の80\,\% が電子を1個ずつ押し出す（量子効率 0.8）とき、光 \qty{1}{W} あたりの電流（感度）を求めましょう。\qty{1}{\micro W} の光では、何 A 流れますか。
\phbamen
光の信号を電流に変えるセンサーの性能を求めます。
\phmushi
光の反射や、電子が電極に着くまでに消えることは、量子効率にふくめて考えます。
\phzu
$\lambda=\qty{850}{nm}$、$\eta=0.8$、$hc=\qty{1240}{eV.nm}$。
\phriyu
光 1 W あたりの光子の数は $1/(hf)$、1個が $\eta$ 個の電子を押し出すので、⑩の感度 $=\eta e/(hf)$ です。
\phtateru
\[ \text{感度}=\eta\frac{e}{hf}=\eta\frac{\lambda\,\text{〔nm〕}}{1240}\ \text{A/W} \]
\phtoku
光子1個のエネルギーは $1240/850\approx\qty{1.46}{eV}$。感度 $=0.8\div1.46\approx\qty{0.55}{A/W}$。\qty{1}{\micro W} では $\qty{0.55}{\micro A}$。
\phtashikame
単位：$e/(hf)$ は C÷J＝A/W。光子のエネルギー $E$ を eV で書くと、$e/(hf)=1/E$〔A/W〕。\\
大きさ：光子1個で電子1個なら、1.46 eV の光子 1 W 分で、いちばん大きい値は $1/1.46\approx\qty{0.68}{A}$ です。\\
極端な場合：波長が 1110 nm（ケイ素の $E_{\mathrm g}$）より長いと、光子は吸われず、感度は0になります。
\kotae 約\qty{0.55}{A/W}、\qty{1}{\micro W} で約\qty{0.55}{\micro A}
\end{phreidai}
\begin{jibun}
波長 \qty{650}{nm}、量子効率 0.7 のフォトダイオードの感度を求めましょう。\qty{2}{mW} の光では何 mA 流れますか。
\bkjans{$0.7\times650\div1240\approx\qty{0.37}{A/W}$。\qty{2}{mW} で約\qty{0.73}{mA}}
\end{jibun}

\label{ph:原-9:ex13}
\begin{phreidai}[ペルチェ素子の冷蔵庫の上限]
ペルチェ素子の小さな冷蔵庫で、中を 5 ℃に保ち、35 ℃の部屋に熱を捨てます。成績係数（くみ出す熱÷使う電気）の上限を求めましょう。
\phbamen
物理シリーズの『熱学』（熱-5）が約束した、ペルチェ素子と第二法則の確かめです。
\phmushi
上限なので、素子の損失（ジュール熱・熱のもどり）はないものとします。
\phzu
$T_{\mathrm L}=278$ K、$T_{\mathrm H}=308$ K。
\phriyu
どんなヒートポンプも、カルノーの成績係数（熱-5 ⑤）をこえられません。ペルチェ素子の中で熱を運ぶのが電子でも、第二法則は変わりません。
\phtateru
\[ \mathrm{COP}_{\text{冷}}\leqq\frac{T_{\mathrm L}}{T_{\mathrm H}-T_{\mathrm L}} \]
\phtoku
\[ \mathrm{COP}_{\text{冷}}\leqq\frac{278}{308-278}=\frac{278}{30}=9.26\cdots \]
上限なので切り捨てて、約 9.2 です。
\phtashikame
単位：K÷K で、単位のない数。\\
熱-5 との照合：熱-5 の章末問題と同じ $278/30$。熱-5 では四捨五入して 9.3 と書きました。\\
実際：ペルチェ素子の成績係数は、ふつう1より小さい。上限との差は、柱の中のジュール熱と、熱い面から冷たい面へもどる熱伝導によるものです。
\kotae 上限は約9.2
\end{phreidai}
\begin{jibun}
ペルチェ素子で、中を 0 ℃、外を 40 ℃にするときの成績係数の上限を求めましょう。
\bkjans{$273\div40\approx6.8$（温度の差が大きいほど、上限は小さくなる）}
\end{jibun}

%% ======================== まちがえやすいところ ========================
\Needspace{11\baselineskip}\bksection{tsumazuki}{まちがえやすいところをチェック}
\begin{tsumazuki}{正孔は、＋のイオン？}
\bkng{p形の正孔は、ホウ素の原子が＋になったもの。}
\bkok{正孔は、結合の手から電子が抜けた穴で、結晶の中を動く。ホウ素は電子を受けとって－のイオンになり、動かない。}
空乏層の中に残るのは、動けないイオンです（⑥）。正孔とイオンを取り違えないようにしましょう。
\end{tsumazuki}

\begin{tsumazuki}{n形の半導体は、－の電気をおびている？}
\bkng{n形は電子が多いので、全体が－になっている。}
\bkok{n形も p形も、全体の電荷は0。電子の数だけ、ドナーのイオンの＋がある。}
n は「動けるキャリアが負（電子）」の意味です（②）。
\end{tsumazuki}

\begin{tsumazuki}{ダイオードは、0.6 V までまったく流れない？}
\bkng{0.6 V より小さい電圧では、電流は0。}
\bkok{電流は、0.06 V ごとに10倍ずつ、指数関数で増え続ける。0.6 V あたりで、mA ほどの目に見える大きさになるだけ。}
片対数グラフ（図\ref{ph:原-9:fig-iv}の右）で見ると、立ち上がりはありません。
\end{tsumazuki}

\begin{tsumazuki}{太陽電池は、つなげばつなぐほど電力が増える？}
\bkng{太陽電池にどんな負荷をつないでも、同じ電力が取り出せる。}
\bkok{電力が最大になる電圧（$V_{\mathrm{oc}}$ より少し低い）は1つ。短絡でも開放でも電力は0。}
例題10 のように、負荷の電圧を選ぶ必要があります。
\end{tsumazuki}

\begin{tsumazuki}{LED の色は、電圧で変わる？}
\bkng{LED に大きな電圧をかけると、青くなる。}
\bkok{色は材料の禁制帯幅 $E_{\mathrm g}$ で決まる。電圧を上げると、電流が増えて明るくなるだけ。}
光子のエネルギーは $hf\approx E_{\mathrm g}$ です（⑨）。
\end{tsumazuki}


%% ======================== 工学ではこう使う ========================
\Needspace{10\baselineskip}\bksection{phgear}{工学ではこう使う}
\begin{phkougaku}[光と電気を行き来する部品]
\textbf{光る}：照明・信号機・スマホの画面の明かり（LED）、光ファイバー通信の光源（半導体レーザー）、テレビのリモコン（赤外 LED）。白い LED は、青い LED と黄色の蛍光体を組み合わせています（⑨）。

\textbf{電気をつくる}：太陽光発電のパネル（⑩）。セルを60〜72枚ほど直列にして 30〜40 V ほどにし、パワーコンディショナー（インバーター：電-16 ⑦）で家の交流に変えます。

\textbf{はかる}：カメラの画素・光ファイバーの受け手（フォトダイオード）、体温計・エアコンの温度センサー（サーミスターやダイオードの電圧）、放射線の検出器（原-14）。

\textbf{整える}：充電器・パソコンの電源の整流（ダイオード：⑦）。電気自動車の充電器や太陽光発電のインバーターでは、ケイ素より禁制帯の大きい炭化ケイ素（SiC）のダイオードも使われます。$n_{\mathrm i}$ が小さいので、高い温度でも逆方向の電流が増えにくいからです。

\begin{phketa}[半導体の数の目安]
\phketaline{熱の電圧 $k_{\mathrm B}T/e$}{\qty{25.9}{mV}（300 K）}
\phketaline{ケイ素のダイオードの順方向}{0.6〜0.7 V、60 mV で電流10倍}
\phketaline{ケイ素の $n_{\mathrm i}$}{約 $\qty{1e10}{cm^{-3}}$、原子 $10^{13}$ 個に1組}
\phketaline{太陽電池のセル1枚}{0.5〜0.7 V、効率 20 \% 前後から 25 \% ほど}
\phketaline{空乏層}{約 \qty{0.1}{\micro m}〜数 \si{\micro m}、電場 $10^4$〜$10^5$ V/cm}
\phboxhead{答えが桁で合っているか}
キャリアの数は、必ず $np=n_{\mathrm i}^{\,2}$ で確かめましょう。n形で $p$ が $n_{\mathrm i}$ より大きくなったら、どこかで割り算を逆にしています。ダイオードの電圧が 1 V を大きくこえたり、0.3 V より小さくなったりしたら、$\ln$ の中の電流の単位（mA と A）を疑いましょう。
\end{phketa}
\end{phkougaku}
\phstation{電子デバイス,太陽光発電,光通信,照明工学}{キャリア、ドナー、アクセプター、移動度、空乏層、内蔵電位、順方向・逆方向、飽和電流、開放電圧、曲線因子、量子効率}

\begin{column}{点接触から pn 接合へ}
ラジオが始まった100年ほど前、鉱石の結晶に細い金属の針を立てたもの（鉱石検波器）が、電波の信号を整流する部品として使われていました。なぜ一方向にだけ流れるのかは、長いあいだ謎でした。1930〜40年代、量子力学で固体のバンドがわかってくると、半導体と金属の境目や、n形と p形の境目の「坂」で説明できるようになりました。1949年、ショックレーは pn 接合の理論をまとめ、⑦のダイオードの式を導きました。トランジスタの発明（原-10）と、この理論により、半導体は、ラジオの検波器から、すべての電子機器の土台になったのです（ショックレー・バーディーン・ブラッテンは、1956年のノーベル物理学賞）。
\end{column}

%% ======================== 章末問題 ========================
\bkendproblems
\bkproblemgroup{基本}{3}
\bkproblem 黄色の LED は、波長 \qty{590}{nm} の光を出します。材料の禁制帯幅は約何 eV ですか。また、禁制帯幅 \qty{3.4}{eV} の窒化ガリウムだけで LED を作ると、何 nm の光が出ますか。目に見えますか。
\bkpans{$1240\div590\approx\qty{2.1}{eV}$。$1240\div3.4\approx\qty{365}{nm}$ で紫外線（目に見えない）}
\bkproblem ケイ素にヒ素（電子を5個もつ原子）を $\qty{5e15}{cm^{-3}}$ 混ぜました。300 K の電子と正孔の数密度、抵抗率を求めましょう（$n_{\mathrm i}=\qty{1.0e10}{cm^{-3}}$、電子の移動度 1350 cm$^2$/(V·s)）。
\bkpans{電子 $\qty{5e15}{cm^{-3}}$、正孔 $10^{20}\div(5\times10^{15})=\qty{2e4}{cm^{-3}}$。$\sigma=1.60\times10^{-19}$ $\times5\times10^{15}\times1350$ $\approx\qty{1.08}{S/cm}$、抵抗率 約\qty{0.93}{\ohm.cm}}
\bkproblem 飽和電流 $\qty{2e-14}{A}$ のダイオードに \qty{5}{mA} を流す電圧を求めましょう。電流を1000倍にするには、電圧をあと何 V 上げればよいですか（300 K）。
\bkpans{$0.0259\times\ln(2.5\times10^{11})\approx\qty{0.68}{V}$。$0.0259\times\ln1000\approx\qty{0.18}{V}$}
\bkproblem 波長 \qty{940}{nm}（リモコンの赤外線）、量子効率 0.6 のフォトダイオードの感度を求めましょう。\qty{0.50}{mW} の光が当たると、何 mA 流れますか。
\bkpans{$0.6\times940\div1240\approx\qty{0.45}{A/W}$。$0.455\times0.50\approx\qty{0.23}{mA}$}
\bkproblem 光電流の密度 \qty{32}{mA/cm^2}、飽和電流の密度 $\qty{1e-12}{A/cm^2}$ のセルの開放電圧を求めましょう（300 K）。このセルを72枚直列にしたパネルの開放電圧はいくらですか。
\bkpans{$0.0259\times\ln(3.2\times10^{10})\approx\qty{0.63}{V}$。72枚で約\qty{45}{V}}
\bkproblem ケイ素の真性キャリア密度は、300 K から 400 K に上げると何倍になりますか。例題2 と同じ $N_{\mathrm c}$・$N_{\mathrm v}$・$E_{\mathrm g}$ を使いましょう（400 K の $k_{\mathrm B}T=\qty{0.03447}{eV}$）。
\bkpans{$(400/300)^{3/2}\times e^{21.66-16.25}=1.54\times224\approx350$ 倍}
\bkproblem ペルチェ素子の冷蔵庫で、中を 7 ℃（280 K）、外を 35 ℃（308 K）にして、1秒に \qty{20}{J} の熱をくみ出します。（1）成績係数の上限と、そのとき要る電力を求めましょう。（2）実際の成績係数が 0.5 なら、要る電力はいくらですか。
\bkpans{（1）\nobreak 上限 $280\div28=10$。電力は少なくとも $20\div10=\qty{2.0}{W}$。（2）\nobreak $20\div0.5=\qty{40}{W}$}
\bkproblem 波長 \qty{630}{nm} の赤い LED に、\qty{2.0}{V} で \qty{20}{mA} を流します。（1）LED が使う電力はいくらですか。（2）その 30\,\% が光になるとき、1秒に出る光子の数を求めましょう（$hc=\qty{1240}{eV.nm}$、$\qty{1}{eV}=\qty{1.60e-19}{J}$）。
\bkpans{（1）\nobreak $2.0\times0.020=\qty{0.040}{W}$。（2）\nobreak 光は $0.30\times0.040=\qty{0.012}{W}$。光子1個は $1240\div630\approx\qty{1.97}{eV}$ $\approx\qty{3.15e-19}{J}$ なので、$0.012\div(3.15\times10^{-19})\approx3.8\times10^{16}$ 個}
\bkproblemgroup{チャレンジ}{5}
\bkproblem アクセプターが $\qty{1e18}{cm^{-3}}$ の p形と、ドナーが $\qty{1e15}{cm^{-3}}$ の n形でできた pn 接合（300 K、$n_{\mathrm i}=\qty{1.0e10}{cm^{-3}}$、比誘電率 11.7）の、内蔵電位・空乏層の厚さ・いちばん強い電場を求めましょう。空乏層の何\,\% が n形の側にありますか。
\bkpans{$V_{\mathrm{bi}}=0.0259\times\ln10^{13}\approx\qty{0.78}{V}$、$W\approx\qty{1.0}{\micro m}$、$|\mathcal E_{\max}|=2V_{\mathrm{bi}}/W\approx\qty{1.5e4}{V/cm}$。$x_{\mathrm n}/W=N_{\mathrm a}/(N_{\mathrm a}+N_{\mathrm d})$ で、約99.9\,\% が n形の側}
\bkproblem p$^+$n 接合のダイオード（面積 \qty{4}{mm^2}、n形の $N_{\mathrm d}=\qty{1e15}{cm^{-3}}$、正孔の移動度 480 cm$^2$/(V·s)、寿命 \qty{10}{\micro s}）の飽和電流を、⑫のショックレーの式で求めましょう（300 K、$n_{\mathrm i}=\qty{1.0e10}{cm^{-3}}$）。順方向に \qty{0.60}{V} かけると、何 mA 流れますか。
\bkpans{$D_{\mathrm p}\approx\qty{12.4}{cm^2/s}$、$L_{\mathrm p}=\sqrt{12.4\times10^{-5}}\approx\qty{0.011}{cm}$。$I_{\mathrm s}=eSn_{\mathrm i}^{\,2}D_{\mathrm p}$ $\div(L_{\mathrm p}N_{\mathrm d})$ $\approx\qty{7.1e-13}{A}$。\qty{0.60}{V} で約\qty{8.2}{mA}}
\bkproblem （1）アクセプターを $\qty{1e17}{cm^{-3}}$ 入れた p形のケイ素で、フェルミ準位は価電子帯の上のはしから何 eV 上にありますか（$N_{\mathrm v}=\qty{1.04e19}{cm^{-3}}$）。（2）ドナーを $\qty{1e18}{cm^{-3}}$ 入れた n形では、伝導帯の下のはしから何 eV 下ですか（$N_{\mathrm c}=\qty{2.8e19}{cm^{-3}}$）。（3）（2）は $k_{\mathrm B}T$ の何倍ですか。ボルツマンの裾の近似は、まだ使えそうですか。
\bkpans{（1）\nobreak $0.0259\times\ln104\approx\qty{0.12}{eV}$。（2）\nobreak $0.0259\times\ln28\approx\qty{0.086}{eV}$。（3）\nobreak 約3.3倍。$e^{3.3}\approx27$ なので $+1$ の無視で4\,\% ほどずれる。まだ使えるが、これより濃くすると近似がくずれる}
\bkproblem ペルチェ素子を発電に使います。1対で 1 K あたり \qty{0.4}{mV} の素子が100対あり、内部の抵抗は \qty{2}{\ohm} です。両面の温度の差を 50 K にしたとき、（1）何もつながないときの電圧、（2）\qty{2}{\ohm} の負荷をつないだときの電流と、負荷で使われる電力を求めましょう。
\bkpans{（1）\nobreak $0.4\times10^{-3}\times100\times50=\qty{2.0}{V}$。（2）\nobreak 電流 $2.0\div(2+2)=\qty{0.50}{A}$、電力 $2\times0.50^2=\qty{0.50}{W}$（内部抵抗と同じ負荷で、いちばん大きな電力：電-5 ⑨）}

\begin{furikaeri}
\begin{checklist}
\item 光電池に光を当てると電気ができ、LED は決まった向きにつないだときだけ光ることを、実験で確かめられる。\bkfurref{ph:原-9:sec1}
\item n形・p形の半導体の電子と正孔の数と抵抗率を求め、pn 接合が電流を一方の向きにだけ通すしくみを、図で説明できる。\bkfurref{ph:原-9:sec4}
\item ダイオードの電流と電圧の関係（約 60 mV で電流10倍）、LED の色と禁制帯の幅の関係を使える。\bkfurref{ph:原-9:sec7}
\item 太陽電池の電流と電圧のグラフから、開放電圧・いちばん大きな電力・曲線因子・効率を求められる。\bkfurref{ph:原-9:sec10}
\item フェルミ分布の裾からキャリアの数の式を、ドリフトと拡散のつり合いから内蔵電位とダイオードの式を導ける。\bkfurref{ph:原-9:sec11}
\end{checklist}
\end{furikaeri}
\bkfukidashi{sora}{LED は、電子が坂を下りるときのエネルギーを、そのまま光にしていたんだね。光電池はその逆で、光が電子を坂の上に持ち上げていたなんて、同じ部品の表と裏みたい！}
